% =====================================================================
% Restoring Generation Individuality from Axiom 0 Alone
% Chapter for the S2-11DM2ET-X Axiomatic Book
% % =====================================================================
% Restoring Generation Individuality from Axiom 0 Alone
% Chapter for the S2-11DM2ET-X Axiomatic Book
% % =====================================================================
% Restoring Generation Individuality from Axiom 0 Alone
% Chapter for the S2-11DM2ET-X Axiomatic Book
% % =====================================================================
% Restoring Generation Individuality from Axiom 0 Alone
% Chapter for the S2-11DM2ET-X Axiomatic Book
% \input{Individuality_from_Axiom0_Expansion}  (after Axiom 0 / Math 1)
% =====================================================================

\chapter{Restoring Generation Individuality from Axiom 0}
\label{ch:individuality}

\section{Motivation and Problem Statement}

\begin{definition}[Multiplicity vs.\ Individuality]
\emph{Multiplicity} is the cardinality of the generation set. \emph{Individuality}
is a complete set of inequivalent labels distinguishing the three Standard-Model
generations as ordered individuals \((g_0,g_1,g_2)\), together with observables
(masses, mixings, lifetimes) that depend nontrivially on those labels.
\end{definition}

\begin{lemma}[Multiplicity from Axiom 0]
\label{lem:multiplicity}
Axiom 0 asserts that the universe contains exactly three generations of
Standard-Model fermions. This forces the qutrit structure, the non-perturbative
superpotential \(W_{np}=e^3\), the flux budget
\[
N_{\mathrm{flux}}=\lfloor e^3\times 3^5\rfloor=4880,
\]
HQCC termination in \(539\pm 1\) steps, the immutable breathing mode
\(G_4=539.9\,\mathrm{s}\), and the puncture count \(|P|=61\).
\end{lemma}

\begin{lemma}[Individuality Gap]
\label{lem:gap}
Axiom 0 alone determines multiplicity but does not determine individuality.
Any two bijections of the generation set are observationally equivalent at the
level of Axiom 0: the axiom is invariant under \(\mathrm{Sym}(\{0,1,2\})\).
\end{lemma}

\begin{corollary}[Need for Additional Discrete Structure]
To restore individuality while still beginning from the single empirical fact of
three generations, additional structure is required. The structure must be
\emph{discrete}, forced by integers already present in the derived chain
\(\{3,61,539,1001,4880,\ldots\}\), and must not reopen the fixed-point rigidity
of \(G_4\) or the puncture count \(|P|=61\).
\end{corollary}

\begin{definition}[Admissible Individuality Structure]
An \emph{admissible individuality structure} is a map
\[
\iota\colon \{g_0,g_1,g_2\}\longrightarrow \mathcal{L}
\]
into a discrete label space \(\mathcal{L}\) such that:
\begin{enumerate}[label=(\roman*)]
  \item \(\iota\) is injective (labels distinguish generations);
  \item \(\mathcal{L}\) is constructed only from Axiom 0 and the Global Lemma Core
        (no free continuous parameters);
  \item the push-forward of \(\iota\) leaves \(G_4=539.9\,\mathrm{s}\) and
        \(|P|=61\) invariant;
  \item residual \(\mathrm{Sym}(\{0,1,2\})\) is broken only by integers already
        fixed in the chain.
\end{enumerate}
\end{definition}

Three admissible routes are developed below and closed by explicit integer
formulae (Section~\ref{sec:closed}). They compose into a single package
(Section~\ref{sec:composite}) with a canonical charged-lepton table
(Section~\ref{sec:table}).

% ---------------------------------------------------------------------
\section{Shared Scaffolding}
\label{sec:scaffold}

\begin{definition}[Balanced-Ternary Syracuse Map]
\label{def:T3}
The balanced-ternary Syracuse map \(T_3\colon\mathbb{N}\to\mathbb{N}\) is the
integer-valued Collatz-type operator of the HQCC theorem (exact \(539\)-step
termination above the model threshold). Write \(\sigma(n)\) for the stopping
time to the fixed point \(1\).
\end{definition}

\begin{definition}[Residual \(+1\) Charge]
\label{def:residual}
The residual charge is the puncture class modulo the generation qutrit:
\begin{equation}
\label{eq:Qres}
Q_{\mathrm{res}}\;:=\; |P|\bmod 3 \;=\; 61\bmod 3 \;=\; 1.
\end{equation}
Equivalently, \(Q_{\mathrm{res}}=+1\in\mathbb{Z}/3\mathbb{Z}\) is the leftover unit
after balanced-ternary cancellation on the bulk cycles. (Note:
\(N_{\mathrm{flux}}=4880\equiv 2\pmod{3}\); the residual used for individuality is
the puncture residual \eqref{eq:Qres}, not the raw flux residue.)
\end{definition}

\begin{definition}[11-Dimensional Bulk and Cycles]
\label{def:11D}
Let \(X^{11}\) be the 11-dimensional bulk of the model, with temporal torsion
operator
\[
T=d+(\zeta^t-1)\wedge H_3,\qquad
\zeta=\exp\bigl(2\pi i/539.9\bigr).
\]
Let \(\mathcal{C}_k(X^{11})\) denote integral \(k\)-cycles relevant to flux
wrapping (in particular those supporting \(G_4\) and the 61 punctures on the
\(G_2\)-holonomy 7-manifold).
\end{definition}

\begin{definition}[1001-Simplex]
\label{def:1001}
Let \(\Delta^{1001}\) be the standard 1001-simplex. The integer
\[
1001=7\times 11\times 13
\]
is already present in the derived numerical chain (e.g.\ the factor \(61/1001\)
in lifetime formulae). Shells of \(\Delta^{1001}\) are discrete orbits under a
specified group action or filtration (Definition~\ref{def:shells}); they are
\emph{not} continuous radial coordinates.
\end{definition}

\begin{lemma}[Fixed-Point Rigidity Constraint]
\label{lem:rigidity}
Any admissible individuality structure must preserve the unique fixed point
\[
G_4=539.9\,\mathrm{s}
\]
satisfying simultaneously HQCC termination, flux quantization, puncture count,
entropy saturation, and the sub-harmonic set \(\{5,10,15,30,45\}\,\mathrm{s}\).
In particular, one may not introduce three independent fundamental periods.
\end{lemma}

\begin{definition}[Generation Index]
\label{def:index}
Write \(j\in\{0,1,2\}\) for the generation index, with the charged-lepton
identification of Section~\ref{sec:table}:
\[
j=0\leftrightarrow e,\qquad
j=1\leftrightarrow \mu,\qquad
j=2\leftrightarrow \tau.
\]
The same index labels the three quark generations \((d,s,b)\) and \((u,c,t)\)
when individuality is lifted from leptons to quarks.
\end{definition}

% ---------------------------------------------------------------------
\section{Route I: Topological Charges and Winding Numbers}
\label{sec:route1}

\begin{definition}[Generation Winding Data]
\label{def:windings}
A \emph{generation winding assignment} is a triple
\[
\mathbf{w}=(w_0,w_1,w_2)\in H_k\bigl(X^{11};\mathbb{Z}\bigr)^3
\]
(or integer winding numbers of three distinguished cycles
\(\gamma_0,\gamma_1,\gamma_2\in\mathcal{C}_k(X^{11})\)) such that
\begin{equation}
\label{eq:winding-sum}
\sum_{j=0}^{2} w_j \;=\; W_{\mathrm{tot}},
\end{equation}
where \(W_{\mathrm{tot}}\) is an integer fixed by the closed formulae of
Section~\ref{sec:closed} (no free continuous modulus).
\end{definition}

\begin{definition}[Arithmetic Progression Form]
\label{def:AP}
The assignment \(\mathbf{w}\) is in \emph{arithmetic progression form} if there
exist integers \(w_0,\Delta\) with
\begin{equation}
\label{eq:AP}
w_j \;=\; w_0 + j\,\Delta,\qquad j\in\{0,1,2\},
\end{equation}
and hence
\begin{equation}
\label{eq:AP-sum}
W_{\mathrm{tot}} \;=\; 3w_0 + 3\Delta \;=\; 3(w_0+\Delta).
\end{equation}
\end{definition}

\begin{theorem}[Route I Individuality]
\label{thm:route1}
Suppose \(\mathbf{w}\) is the canonical arithmetic-progression assignment of
Section~\ref{sec:closed}. Then
\[
\iota_I\colon g_j\mapsto w_j
\]
is an admissible individuality structure. Generations are distinguished by
inequivalent topological charges on the 11-dimensional cycles, while
multiplicity and global flux remain Axiom-0 data.
\end{theorem}

\begin{proof}[Proof sketch]
Injectivity follows from \(\Delta=61\neq 0\). Discreteness is immediate from
integral homology. Invariance of \(G_4\) and \(|P|\) follows because
\eqref{eq:winding-sum} preserves the total winding class used for the breathing
mode and puncture count; only the distribution among three cycles changes.
Residual \(\mathrm{Sym}_3\) is broken by the ordered progression
\(w_0<w_1<w_2\), with orientation fixed by the chirality of \(T\).
\end{proof}

\begin{corollary}[Mass Hierarchy from Wrapping]
If particle masses scale with wrapping volume (or with the period of \(T\) along
\(\gamma_j\)), then
\[
m_j \;\propto\; \mathrm{Vol}(\gamma_j)
\quad\text{or}\quad
m_j \;\propto\; \bigl|\langle T,[ \gamma_j]\rangle\bigr|
\]
yields a discrete hierarchy indexed by \(w_j\), with no free continuous Yukawa
modulus beyond integers already in the chain.
\end{corollary}

\begin{remark}[Failure Mode]
If \(\mathbf{w}\) is chosen from an open family of windings not constrained by
\eqref{eq:winding-sum} and Section~\ref{sec:closed}, individuality is smuggled
in as a second axiom. Route I is admissible only in the forced form.
\end{remark}

% ---------------------------------------------------------------------
\section{Route II: Leftover \(+1\) on Shells of the 1001-Simplex}
\label{sec:route2}

\begin{definition}[Shells of \(\Delta^{1001}\)]
\label{def:shells}
Let \(\Gamma=\mathbb{Z}_7\times\mathbb{Z}_{11}\times\mathbb{Z}_{13}\) act by the
prime factorization \(1001=7\times 11\times 13\) (equivalently,
\(\Gamma_3=\mathbb{Z}_3\) by the generation qutrit). A \emph{shell} is a
\(\Gamma\)-orbit or a grade of a \(\Gamma\)-equivariant filtration
\[
\Delta^{1001}\;=\;\bigsqcup_{\ell} S_\ell.
\]
Shell indices are integers; there is no continuous radial coordinate.
\end{definition}

\begin{definition}[Generation-Dependent Embedding of \(Q_{\mathrm{res}}\)]
\label{def:embed}
A \emph{residual embedding} assigns the leftover charge \(Q_{\mathrm{res}}=+1\)
to three inequivalent shells:
\begin{equation}
\label{eq:shells}
\iota_{II}\colon g_j \mapsto S_{k_j},\qquad
k_0,k_1,k_2\text{ pairwise distinct}.
\end{equation}
\end{definition}

\begin{lemma}[Single-Shell Obstruction]
\label{lem:single-shell}
If \(Q_{\mathrm{res}}=+1\) is placed in a single shell only, then either
\begin{enumerate}[label=(\alph*)]
  \item the three-generation qutrit collapses to a single effective grade
        (loss of multiplicity), or
  \item flux quantization / entropy conservation across the \(+U/-U\) interface
        fails to match the Global Lemma Core.
\end{enumerate}
Hence any residual embedding compatible with Axiom 0 must occupy (at least)
three inequivalent shells.
\end{lemma}

\begin{theorem}[Route II Individuality]
\label{thm:route2}
Let \(\iota_{II}\) be the canonical residual embedding of
Section~\ref{sec:closed} with shell indices \(\{7,11,13\}\). Then \(\iota_{II}\)
is an admissible individuality structure.
\end{theorem}

\begin{proof}[Proof sketch]
Injectivity is pairwise distinct primes. Discreteness is by
Definition~\ref{def:shells}. The total residual remains \(+1\), so \(G_4\) and
\(|P|\) are unchanged. The identity \(7\times 11\times 13=1001\) and the factor
\(61/1001\) already appear in lifetime formulae.
\end{proof}

\begin{corollary}[Lifetimeage Coupling]
The model lifetime schematic
\[
\tau \;\propto\; \frac{61}{1001}\cdot\frac{\phi^7}{539.9}
\]
admits the generation refinement
\begin{equation}
\label{eq:tau-g}
\tau_j \;\propto\; \frac{61}{1001}\cdot f(k_j)\cdot\frac{\phi^7}{539.9},
\end{equation}
with the fixed rational weight of Section~\ref{sec:closed}
(no free continuous coupling).
\end{corollary}

\begin{remark}[Failure Mode]
Continuous ``shell radius'' reintroduces parameters and is forbidden.
\end{remark}

% ---------------------------------------------------------------------
\section{Route III: Graded Balanced-Ternary Map}
\label{sec:route3}

\begin{definition}[Graded Ternary Data]
\label{def:graded}
A \emph{grading} of the balanced-ternary map is a family of cocycles
\[
c_j\colon\mathbb{N}\to\{-1,0,+1\},\qquad j\in\{0,1,2\},
\]
together with the common underlying map \(T_3\), such that
\begin{equation}
\label{eq:graded-step}
T_3^{(j)}(n)\;:=\;T_3(n)
\quad\text{(set-map unchanged)},
\end{equation}
while observables use weights \(c_j\). The cocycles satisfy
\begin{equation}
\label{eq:cocycle-balance}
\sum_{j=0}^{2} c_j(n)\;=\;Q_{\mathrm{res}}\;=\;+1
\quad\text{for all \(n\) in the master tower}.
\end{equation}
\end{definition}

\begin{definition}[Contraction Rates and Resonant Periods]
\label{def:rates}
Write \(\lambda=\ln 3/539\) for the Banach contraction rate of the ungraded HQCC
map. Graded contractions \(\lambda_j\) must be rational multiples of \(\lambda\).
Graded resonant sub-periods \(\tau_j^{\mathrm{sub}}\) must be drawn from the
sub-harmonic set \(\{5,10,15,30,45\}\,\mathrm{s}\) (or integer commensurators of
\(G_4\))---never three independent fundamental periods.
\end{definition}

\begin{lemma}[Common Orbit Constraint]
\label{lem:common-orbit}
If the three generations are assigned three independent fundamental periods,
fixed-point rigidity (Lemma~\ref{lem:rigidity}) fails. Therefore any admissible
grading keeps a common termination depth \(\sigma=539\pm 1\) and a common
fundamental mode \(G_4\), with individuality in cocycle weights, residue classes
modulo 3, or sub-period structure.
\end{lemma}

\begin{theorem}[Route III Individuality]
\label{thm:route3}
Let \((c_0,c_1,c_2)\) be the canonical graded cocycles of
Section~\ref{sec:closed}. Then
\[
\iota_{III}\colon g_j\mapsto c_j
\]
is an admissible individuality structure: one orbit, one \(G_4\),
generation-dependent dynamical observables.
\end{theorem}

\begin{proof}[Proof sketch]
The set-map \(T_3\) and stopping time \(\sigma\) are generation-independent.
Inequivalence of the \(c_j\) (distinct mod-3 supports) breaks \(\mathrm{Sym}_3\).
Observables (mean sojourn near punctures, weighted transfer spectrum) depend on
\(j\) only through \(c_j\). All numerical inputs are chain integers.
\end{proof}

\begin{remark}[Failure Mode]
Independent fundamental periods not commensurate with \(G_4\) are inadmissible.
\end{remark}

% ---------------------------------------------------------------------
\section{Closed Integer Formulae}
\label{sec:closed}

All ``forced by the chain'' steps are fixed here by explicit integers from
\(\{3,7,11,13,61,539,1001,4880\}\). No continuous choice remains.

\begin{definition}[Canonical Step and Base Winding]
\label{def:Delta}
\[
\Delta \;:=\; |P| \;=\; 61,
\qquad
w_0 \;:=\; \sigma_{\mathrm{HQCC}} \;=\; 539.
\]
\end{definition}

\begin{lemma}[Canonical Windings]
\label{lem:canonical-w}
The unique arithmetic-progression winding assignment with base \(w_0=539\) and
step \(\Delta=61\) is
\begin{equation}
\label{eq:w-explicit}
w_j \;=\; 539 + 61\,j,\qquad j\in\{0,1,2\},
\end{equation}
i.e.
\[
(w_0,w_1,w_2) \;=\; (539,\,600,\,661).
\]
The total class is
\begin{equation}
\label{eq:Wtot}
W_{\mathrm{tot}} \;=\; 539+600+661 \;=\; 1800 \;=\; 3\times 600 \;=\; 3(w_0+\Delta).
\end{equation}
\textbf{Status (No-Go chapter).}
The pair \((w_0,\Delta)=(539,61)\) is \emph{conditional} on a non-circular supply
of \(\sigma=539\) (and of \(|P|=61\)). Flux democracy alone does not select
these anchors (Theorem on flux democracy no-go). By contrast, cocycles
\(c_j\) and shells \(k_j\in\{7,11,13\}\) remain a priori.
\end{lemma}

\begin{proof}
Substitute \eqref{eq:w-explicit} into \eqref{eq:AP-sum}. Summation is elementary.
Both \(539\) and \(61\) are Global Lemma Core integers; no other pair is used.
\end{proof}

\begin{lemma}[Canonical Shells]
\label{lem:canonical-k}
The unique unordered triple of prime shells forced by \(1001=7\times 11\times 13\)
is \(\{7,11,13\}\). Ordering by increasing prime (aligned with increasing
\(j\) and chirality of \(T\)) gives
\begin{equation}
\label{eq:k-explicit}
(k_0,k_1,k_2) \;=\; (7,\,11,\,13).
\end{equation}
\end{lemma}

\begin{definition}[Canonical Dictionaries \(\Phi,\Psi\)]
\label{def:PhiPsi}
\begin{align}
\Phi(w_j) &\;:=\; k_j \;=\; p_{j+1},
  \label{eq:Phi}\\
\Psi(k_j) &\;:=\; c_j,
  \label{eq:Psi}
\end{align}
where \(p_1=7\), \(p_2=11\), \(p_3=13\) are the prime factors of \(1001\) in
increasing order, and \(c_j\) is Definition~\ref{def:c-explicit}. Equivalently,
\[
j \;=\; \frac{w_j-539}{61} \;=\; \mathrm{index}(k_j\text{ in }(7,11,13)),
\]
so \(\Phi\) and \(\Psi\) are determined by the single generation index \(j\).
\end{definition}

\begin{definition}[Canonical Cocycles]
\label{def:c-explicit}
\begin{equation}
\label{eq:c-explicit}
c_j(n) \;:=\;
\begin{cases}
+1 & \text{if } n\equiv j \pmod{3},\\
0 & \text{otherwise}.
\end{cases}
\end{equation}
\end{definition}

\begin{lemma}[Cocycle Balance]
\label{lem:cocycle-ok}
For every \(n\in\mathbb{N}\),
\[
\sum_{j=0}^{2} c_j(n) \;=\; 1 \;=\; Q_{\mathrm{res}}.
\]
Thus \eqref{eq:cocycle-balance} holds pointwise on the master tower.
\end{lemma}

\begin{proof}
Exactly one residue class modulo 3 equals \(j\) for a unique \(j\in\{0,1,2\}\).
\end{proof}

\begin{definition}[Canonical Sub-Harmonics and Contractions]
\label{def:sub-explicit}
\begin{equation}
\label{eq:sub-explicit}
\bigl(\tau_0^{\mathrm{sub}},\tau_1^{\mathrm{sub}},\tau_2^{\mathrm{sub}}\bigr)
\;=\;
(5,\,15,\,45)\,\mathrm{s}
\;\subset\;
\{5,10,15,30,45\}\,\mathrm{s},
\end{equation}
ordered lightest \(\leftrightarrow\) shortest, and
\begin{equation}
\label{eq:lambda-explicit}
\lambda_j \;=\; \lambda \;=\; \frac{\ln 3}{539}
\quad\text{for all }j
\end{equation}
(common contraction; generation dependence lives in \(c_j\) and
\(\tau_j^{\mathrm{sub}}\) only).
\end{definition}

\begin{definition}[Canonical Lifetime Weight]
\label{def:f}
\begin{equation}
\label{eq:f}
f(k) \;:=\; \frac{1001}{k},
\qquad
\bigl(f(7),f(11),f(13)\bigr) \;=\; (143,\,91,\,77).
\end{equation}
\end{definition}

\begin{theorem}[Closure of All Forced Steps]
\label{thm:closure}
With Definitions~\ref{def:Delta}--\ref{def:f} and
Lemmas~\ref{lem:canonical-w}--\ref{lem:cocycle-ok}, every open ``forced by the
chain'' step in Routes I--III is an explicit integer formula. The only free
discrete choice retained is the global orientation
(increasing vs.\ decreasing \(j\)), fixed by the chirality of \(T\) to the
increasing convention \(j=0,1,2\leftrightarrow e,\mu,\tau\).
\end{theorem}

% ---------------------------------------------------------------------
\section{Composite Package}
\label{sec:composite}

\begin{definition}[Composite Individuality Structure]
\label{def:composite}
The \emph{composite package} is
\[
\iota \;=\; (\iota_I,\iota_{II},\iota_{III}),
\]
with
\begin{enumerate}[label=(\roman*)]
  \item Route I: \(g_j\mapsto w_j=539+61j\);
  \item Route II: \(w_j\mapsto k_j\in\{7,11,13\}\) via \(\Phi\);
  \item Route III: \(k_j\mapsto c_j\) via \(\Psi\), on the common \(T_3\) orbit.
\end{enumerate}
Consistency is the commuting square
\begin{equation}
\label{eq:dictionary}
\begin{array}{ccc}
\{g_j\} & \xrightarrow{\;\iota_I\;} & \{w_j\} \\[4pt]
\big\downarrow{\scriptstyle \iota_{III}} & & \big\downarrow{\scriptstyle \Phi} \\[4pt]
\{c_j\} & \xleftarrow{\;\Psi\;} & \{k_j\}
\end{array}
\end{equation}
with \(\Phi,\Psi\) as in Definition~\ref{def:PhiPsi}.
\end{definition}

\begin{theorem}[Individuality without a Second Axiom]
\label{thm:main}
The composite package \(\iota\) is an admissible individuality structure:
\begin{enumerate}[label=(\roman*)]
  \item Axiom 0 remains the sole primitive (cardinality three);
  \item the three generations are pairwise inequivalent as labeled individuals;
  \item \(G_4=539.9\,\mathrm{s}\), HQCC depth \(539\pm 1\), and \(|P|=61\) are
        unchanged;
  \item no free continuous parameter is introduced.
\end{enumerate}
\end{theorem}

\begin{proof}[Proof sketch]
Theorem~\ref{thm:closure} supplies explicit integer labels. Routes I--III are
admissible by Theorems~\ref{thm:route1}--\ref{thm:route3}. The dictionaries
\(\Phi,\Psi\) are bijections on a three-element set, hence introduce no continuum
of moduli. Conditions (i)--(iv) follow.
\end{proof}

\begin{corollary}[Symmetry Breaking Pattern]
\[
\mathrm{Sym}(\{0,1,2\})\;\longrightarrow\;\{1\},
\]
broken completely by \(\iota\). Partial activation of a single route may leave
a \(\mathbb{Z}_2\) stabilizer.
\end{corollary}

% ---------------------------------------------------------------------
\section{Canonical Generation-Label Table (\(e/\mu/\tau\))}
\label{sec:table}

\begin{definition}[Canonical Charged-Lepton Dictionary]
\label{def:leptons}
Identify
\[
(g_0,g_1,g_2) \;=\; (e,\,\mu,\,\tau)
\]
with the closed labels of Section~\ref{sec:closed}. The same \(j\) labels quark
generations when lifted: \((d,s,b)\) and \((u,c,t)\).
\end{definition}

\begin{table}[ht]
\centering
\textbf{Table: Canonical individuality labels for charged leptons}
\label{tab:labels}\\[0.5em]
\begin{tabular}{@{}llrrrrl@{}}
\toprule
Gen. & \(j\) & \(w_j\) & \(k_j\) & \(f(k_j)\) & \(\tau_j^{\mathrm{sub}}\) (s) & \(c_j\) support \\
\midrule
\(e\)   & 0 & 539 & 7  & 143 & 5  & \(n\equiv 0\pmod{3}\) \\
\(\mu\) & 1 & 600 & 11 & 91  & 15 & \(n\equiv 1\pmod{3}\) \\
\(\tau\)& 2 & 661 & 13 & 77  & 45 & \(n\equiv 2\pmod{3}\) \\
\bottomrule
\end{tabular}
\end{table}

\begin{table}[ht]
\centering
\textbf{Table: Global invariants preserved by the dictionary}
\label{tab:invariants}\\[0.5em]
\begin{tabular}{@{}ll@{}}
\toprule
Invariant & Value \\
\midrule
\(W_{\mathrm{tot}}=\sum_j w_j\) & \(1800=3\times 600\) \\
\(\prod_j k_j\) & \(7\times 11\times 13=1001\) \\
\(\sum_j c_j(n)\) & \(1=Q_{\mathrm{res}}=61\bmod 3\) \\
\(\lambda_j\) & \(\ln 3/539\) (common) \\
HQCC depth & \(539\pm 1\) \\
\(G_4\) & \(539.9\,\mathrm{s}\) \\
\(|P|\) & \(61\) \\
\bottomrule
\end{tabular}
\end{table}

\begin{corollary}[Explicit Lifetime Refinement]
Under \eqref{eq:tau-g} and \eqref{eq:f},
\begin{equation}
\label{eq:tau-explicit}
\tau_j \;\propto\; \frac{61}{1001}\cdot\frac{1001}{k_j}\cdot\frac{\phi^7}{539.9}
\;=\; \frac{61}{k_j}\cdot\frac{\phi^7}{539.9},
\end{equation}
so the shell weight cancels one factor of \(1001\) and leaves a pure prime in the
denominator: \(\tau_e\propto 61/7\), \(\tau_\mu\propto 61/11\),
\(\tau_\tau\propto 61/13\), times the common \(\phi^7/539.9\) factor. The standard
muon formula in the theorem chain corresponds to the \(j=1\) line with an overall
normalization \(\tau^-\) and resonance correction \(\delta_{\mathrm{res}}\).
\end{corollary}

\begin{corollary}[Discrete Mass Ordering]
Along Route I, winding gaps are constant:
\[
w_1-w_0 \;=\; w_2-w_1 \;=\; 61,
\]
so any observable monotone in \(w_j\) automatically orders
\(e\prec\mu\prec\tau\) (or the reverse if orientation is flipped---excluded by
chirality of \(T\)).
\end{corollary}

\begin{remark}[Quark Lift]
The same table applies to quarks by reinterpreting \(g_j\) as the \(j\)-th quark
generation. Flavour mixing is then the overlap of cocycles \(c_i,c_j\) on the
common HQCC orbit (CKM / PMNS as discrete cocycle correlators), with no
continuous Yukawa matrix beyond the finite dictionary.
\end{remark}

% ---------------------------------------------------------------------
\section{Comparison and Selection}
\label{sec:compare}

\begin{center}
\begin{tabular}{@{}llll@{}}
\toprule
Route & Carrier & Principal risk & Cleanest fit \\
\midrule
I & Windings on 11D cycles & Free winding family & Brane / flux language \\
II & Shells of \(\Delta^{1001}\) & Continuous radius & \(61/1001\) numerics \\
III & Graded \(T_3\) cocycles & Multiple fundamentals & HQCC core \\
Composite & \((w_j,k_j,c_j)\) & Over-constraining \(\Phi,\Psi\) & Full tower \\
\bottomrule
\end{tabular}
\end{center}

\begin{lemma}[Selection under Rigidity]
\label{lem:selection}
If a single route is preferred, Route III with common orbit and graded cocycle
is the least likely to reopen the fixed point of \(G_4\). If full geometric
language is required, the composite package is preferred, with Route III as the
dynamical readout and Routes I--II as geometric realization. The Axiomatic Book
takes the composite package with Table~\ref{tab:labels} as primary.
\end{lemma}

% ---------------------------------------------------------------------
\section{Falsifiability and Downstream Couplings}
\label{sec:falsify}

\begin{definition}[Individuality Observables]
Generation-dependent observables induced by \(\iota\) include:
\begin{enumerate}[label=(\roman*)]
  \item mass ratios \(m_i/m_j\) as functions of \((w_i-w_j)=61(i-j)\);
  \item lifetime ratios from \eqref{eq:tau-explicit}:
        \(\tau_i/\tau_j = k_j/k_i\) at fixed overall normalization;
  \item mixing angles as overlap of cocycles \(c_i,c_j\) on the common orbit;
  \item sub-harmonic phase offsets at \(\tau_j^{\mathrm{sub}}\in\{5,15,45\}\,\mathrm{s}\).
\end{enumerate}
\end{definition}

\begin{lemma}[Falsifiability]
\label{lem:falsify}
The canonical dictionary is falsifiable if measured mass or lifetime ratios
cannot be matched by Table~\ref{tab:labels} (or its single orientation reverse,
already excluded by chirality) while preserving \(G_4\) and \(|P|=61\). There is
no continuous retuning: the label set is finite.
\end{lemma}

\begin{corollary}[Lifetimeage Ratio Test]
At fixed normalization, \eqref{eq:tau-explicit} predicts
\[
\frac{\tau_\mu}{\tau_e} \;=\; \frac{7}{11},\qquad
\frac{\tau_\tau}{\tau_\mu} \;=\; \frac{11}{13},\qquad
\frac{\tau_\tau}{\tau_e} \;=\; \frac{7}{13},
\]
up to the known overall factors already present in the muon theorem line and
channel-dependent phase space. Gross violation of these discrete ratios (after
those factors) falsifies the canonical shell assignment.
\end{corollary}

% ---------------------------------------------------------------------
\section{Global Lemma Core Form}
\label{sec:glc}

\begin{lemma}[Individuality without a Second Axiom]
\label{lem:glc-individuality}
There exists a discrete, Axiom-0-forced assignment
\[
j\;\mapsto\; (w_j,k_j,c_j)
\;=\;
\bigl(539+61j,\; p_{j+1},\; \mathbf{1}_{n\equiv j\pmod{3}}\bigr)
\]
with \((p_1,p_2,p_3)=(7,11,13)\) such that the three generations are inequivalent
while \(G_4\), HQCC depth, and \(|P|\) remain unique.
\end{lemma}

% ---------------------------------------------------------------------
\section{Summary}
\label{sec:summary}

Axiom 0 yields three as a cardinality. Individuality is restored by the
composite integer dictionary:
\begin{enumerate}
  \item windings \(w_j=539+61j\) on 11-dimensional cycles;
  \item leftover \(+1\) on shells \(k_j\in\{7,11,13\}\) of \(\Delta^{1001}\);
  \item graded cocycles \(c_j=\mathbf{1}_{n\equiv j\pmod{3}}\) on a common \(T_3\)
        orbit of length \(539\pm 1\), with sub-harmonics \((5,15,45)\,\mathrm{s}\).
\end{enumerate}
For charged leptons, \(j=0,1,2\leftrightarrow e,\mu,\tau\)
(Table~\ref{tab:labels}). No second continuous axiom is required; fixed-point
rigidity is preserved; falsifiability is finite and discrete.

% end of chapter
  (after Axiom 0 / Math 1)
% =====================================================================

\chapter{Restoring Generation Individuality from Axiom 0}
\label{ch:individuality}

\section{Motivation and Problem Statement}

\begin{definition}[Multiplicity vs.\ Individuality]
\emph{Multiplicity} is the cardinality of the generation set. \emph{Individuality}
is a complete set of inequivalent labels distinguishing the three Standard-Model
generations as ordered individuals \((g_0,g_1,g_2)\), together with observables
(masses, mixings, lifetimes) that depend nontrivially on those labels.
\end{definition}

\begin{lemma}[Multiplicity from Axiom 0]
\label{lem:multiplicity}
Axiom 0 asserts that the universe contains exactly three generations of
Standard-Model fermions. This forces the qutrit structure, the non-perturbative
superpotential \(W_{np}=e^3\), the flux budget
\[
N_{\mathrm{flux}}=\lfloor e^3\times 3^5\rfloor=4880,
\]
HQCC termination in \(539\pm 1\) steps, the immutable breathing mode
\(G_4=539.9\,\mathrm{s}\), and the puncture count \(|P|=61\).
\end{lemma}

\begin{lemma}[Individuality Gap]
\label{lem:gap}
Axiom 0 alone determines multiplicity but does not determine individuality.
Any two bijections of the generation set are observationally equivalent at the
level of Axiom 0: the axiom is invariant under \(\mathrm{Sym}(\{0,1,2\})\).
\end{lemma}

\begin{corollary}[Need for Additional Discrete Structure]
To restore individuality while still beginning from the single empirical fact of
three generations, additional structure is required. The structure must be
\emph{discrete}, forced by integers already present in the derived chain
\(\{3,61,539,1001,4880,\ldots\}\), and must not reopen the fixed-point rigidity
of \(G_4\) or the puncture count \(|P|=61\).
\end{corollary}

\begin{definition}[Admissible Individuality Structure]
An \emph{admissible individuality structure} is a map
\[
\iota\colon \{g_0,g_1,g_2\}\longrightarrow \mathcal{L}
\]
into a discrete label space \(\mathcal{L}\) such that:
\begin{enumerate}[label=(\roman*)]
  \item \(\iota\) is injective (labels distinguish generations);
  \item \(\mathcal{L}\) is constructed only from Axiom 0 and the Global Lemma Core
        (no free continuous parameters);
  \item the push-forward of \(\iota\) leaves \(G_4=539.9\,\mathrm{s}\) and
        \(|P|=61\) invariant;
  \item residual \(\mathrm{Sym}(\{0,1,2\})\) is broken only by integers already
        fixed in the chain.
\end{enumerate}
\end{definition}

Three admissible routes are developed below and closed by explicit integer
formulae (Section~\ref{sec:closed}). They compose into a single package
(Section~\ref{sec:composite}) with a canonical charged-lepton table
(Section~\ref{sec:table}).

% ---------------------------------------------------------------------
\section{Shared Scaffolding}
\label{sec:scaffold}

\begin{definition}[Balanced-Ternary Syracuse Map]
\label{def:T3}
The balanced-ternary Syracuse map \(T_3\colon\mathbb{N}\to\mathbb{N}\) is the
integer-valued Collatz-type operator of the HQCC theorem (exact \(539\)-step
termination above the model threshold). Write \(\sigma(n)\) for the stopping
time to the fixed point \(1\).
\end{definition}

\begin{definition}[Residual \(+1\) Charge]
\label{def:residual}
The residual charge is the puncture class modulo the generation qutrit:
\begin{equation}
\label{eq:Qres}
Q_{\mathrm{res}}\;:=\; |P|\bmod 3 \;=\; 61\bmod 3 \;=\; 1.
\end{equation}
Equivalently, \(Q_{\mathrm{res}}=+1\in\mathbb{Z}/3\mathbb{Z}\) is the leftover unit
after balanced-ternary cancellation on the bulk cycles. (Note:
\(N_{\mathrm{flux}}=4880\equiv 2\pmod{3}\); the residual used for individuality is
the puncture residual \eqref{eq:Qres}, not the raw flux residue.)
\end{definition}

\begin{definition}[11-Dimensional Bulk and Cycles]
\label{def:11D}
Let \(X^{11}\) be the 11-dimensional bulk of the model, with temporal torsion
operator
\[
T=d+(\zeta^t-1)\wedge H_3,\qquad
\zeta=\exp\bigl(2\pi i/539.9\bigr).
\]
Let \(\mathcal{C}_k(X^{11})\) denote integral \(k\)-cycles relevant to flux
wrapping (in particular those supporting \(G_4\) and the 61 punctures on the
\(G_2\)-holonomy 7-manifold).
\end{definition}

\begin{definition}[1001-Simplex]
\label{def:1001}
Let \(\Delta^{1001}\) be the standard 1001-simplex. The integer
\[
1001=7\times 11\times 13
\]
is already present in the derived numerical chain (e.g.\ the factor \(61/1001\)
in lifetime formulae). Shells of \(\Delta^{1001}\) are discrete orbits under a
specified group action or filtration (Definition~\ref{def:shells}); they are
\emph{not} continuous radial coordinates.
\end{definition}

\begin{lemma}[Fixed-Point Rigidity Constraint]
\label{lem:rigidity}
Any admissible individuality structure must preserve the unique fixed point
\[
G_4=539.9\,\mathrm{s}
\]
satisfying simultaneously HQCC termination, flux quantization, puncture count,
entropy saturation, and the sub-harmonic set \(\{5,10,15,30,45\}\,\mathrm{s}\).
In particular, one may not introduce three independent fundamental periods.
\end{lemma}

\begin{definition}[Generation Index]
\label{def:index}
Write \(j\in\{0,1,2\}\) for the generation index, with the charged-lepton
identification of Section~\ref{sec:table}:
\[
j=0\leftrightarrow e,\qquad
j=1\leftrightarrow \mu,\qquad
j=2\leftrightarrow \tau.
\]
The same index labels the three quark generations \((d,s,b)\) and \((u,c,t)\)
when individuality is lifted from leptons to quarks.
\end{definition}

% ---------------------------------------------------------------------
\section{Route I: Topological Charges and Winding Numbers}
\label{sec:route1}

\begin{definition}[Generation Winding Data]
\label{def:windings}
A \emph{generation winding assignment} is a triple
\[
\mathbf{w}=(w_0,w_1,w_2)\in H_k\bigl(X^{11};\mathbb{Z}\bigr)^3
\]
(or integer winding numbers of three distinguished cycles
\(\gamma_0,\gamma_1,\gamma_2\in\mathcal{C}_k(X^{11})\)) such that
\begin{equation}
\label{eq:winding-sum}
\sum_{j=0}^{2} w_j \;=\; W_{\mathrm{tot}},
\end{equation}
where \(W_{\mathrm{tot}}\) is an integer fixed by the closed formulae of
Section~\ref{sec:closed} (no free continuous modulus).
\end{definition}

\begin{definition}[Arithmetic Progression Form]
\label{def:AP}
The assignment \(\mathbf{w}\) is in \emph{arithmetic progression form} if there
exist integers \(w_0,\Delta\) with
\begin{equation}
\label{eq:AP}
w_j \;=\; w_0 + j\,\Delta,\qquad j\in\{0,1,2\},
\end{equation}
and hence
\begin{equation}
\label{eq:AP-sum}
W_{\mathrm{tot}} \;=\; 3w_0 + 3\Delta \;=\; 3(w_0+\Delta).
\end{equation}
\end{definition}

\begin{theorem}[Route I Individuality]
\label{thm:route1}
Suppose \(\mathbf{w}\) is the canonical arithmetic-progression assignment of
Section~\ref{sec:closed}. Then
\[
\iota_I\colon g_j\mapsto w_j
\]
is an admissible individuality structure. Generations are distinguished by
inequivalent topological charges on the 11-dimensional cycles, while
multiplicity and global flux remain Axiom-0 data.
\end{theorem}

\begin{proof}[Proof sketch]
Injectivity follows from \(\Delta=61\neq 0\). Discreteness is immediate from
integral homology. Invariance of \(G_4\) and \(|P|\) follows because
\eqref{eq:winding-sum} preserves the total winding class used for the breathing
mode and puncture count; only the distribution among three cycles changes.
Residual \(\mathrm{Sym}_3\) is broken by the ordered progression
\(w_0<w_1<w_2\), with orientation fixed by the chirality of \(T\).
\end{proof}

\begin{corollary}[Mass Hierarchy from Wrapping]
If particle masses scale with wrapping volume (or with the period of \(T\) along
\(\gamma_j\)), then
\[
m_j \;\propto\; \mathrm{Vol}(\gamma_j)
\quad\text{or}\quad
m_j \;\propto\; \bigl|\langle T,[ \gamma_j]\rangle\bigr|
\]
yields a discrete hierarchy indexed by \(w_j\), with no free continuous Yukawa
modulus beyond integers already in the chain.
\end{corollary}

\begin{remark}[Failure Mode]
If \(\mathbf{w}\) is chosen from an open family of windings not constrained by
\eqref{eq:winding-sum} and Section~\ref{sec:closed}, individuality is smuggled
in as a second axiom. Route I is admissible only in the forced form.
\end{remark}

% ---------------------------------------------------------------------
\section{Route II: Leftover \(+1\) on Shells of the 1001-Simplex}
\label{sec:route2}

\begin{definition}[Shells of \(\Delta^{1001}\)]
\label{def:shells}
Let \(\Gamma=\mathbb{Z}_7\times\mathbb{Z}_{11}\times\mathbb{Z}_{13}\) act by the
prime factorization \(1001=7\times 11\times 13\) (equivalently,
\(\Gamma_3=\mathbb{Z}_3\) by the generation qutrit). A \emph{shell} is a
\(\Gamma\)-orbit or a grade of a \(\Gamma\)-equivariant filtration
\[
\Delta^{1001}\;=\;\bigsqcup_{\ell} S_\ell.
\]
Shell indices are integers; there is no continuous radial coordinate.
\end{definition}

\begin{definition}[Generation-Dependent Embedding of \(Q_{\mathrm{res}}\)]
\label{def:embed}
A \emph{residual embedding} assigns the leftover charge \(Q_{\mathrm{res}}=+1\)
to three inequivalent shells:
\begin{equation}
\label{eq:shells}
\iota_{II}\colon g_j \mapsto S_{k_j},\qquad
k_0,k_1,k_2\text{ pairwise distinct}.
\end{equation}
\end{definition}

\begin{lemma}[Single-Shell Obstruction]
\label{lem:single-shell}
If \(Q_{\mathrm{res}}=+1\) is placed in a single shell only, then either
\begin{enumerate}[label=(\alph*)]
  \item the three-generation qutrit collapses to a single effective grade
        (loss of multiplicity), or
  \item flux quantization / entropy conservation across the \(+U/-U\) interface
        fails to match the Global Lemma Core.
\end{enumerate}
Hence any residual embedding compatible with Axiom 0 must occupy (at least)
three inequivalent shells.
\end{lemma}

\begin{theorem}[Route II Individuality]
\label{thm:route2}
Let \(\iota_{II}\) be the canonical residual embedding of
Section~\ref{sec:closed} with shell indices \(\{7,11,13\}\). Then \(\iota_{II}\)
is an admissible individuality structure.
\end{theorem}

\begin{proof}[Proof sketch]
Injectivity is pairwise distinct primes. Discreteness is by
Definition~\ref{def:shells}. The total residual remains \(+1\), so \(G_4\) and
\(|P|\) are unchanged. The identity \(7\times 11\times 13=1001\) and the factor
\(61/1001\) already appear in lifetime formulae.
\end{proof}

\begin{corollary}[Lifetimeage Coupling]
The model lifetime schematic
\[
\tau \;\propto\; \frac{61}{1001}\cdot\frac{\phi^7}{539.9}
\]
admits the generation refinement
\begin{equation}
\label{eq:tau-g}
\tau_j \;\propto\; \frac{61}{1001}\cdot f(k_j)\cdot\frac{\phi^7}{539.9},
\end{equation}
with the fixed rational weight of Section~\ref{sec:closed}
(no free continuous coupling).
\end{corollary}

\begin{remark}[Failure Mode]
Continuous ``shell radius'' reintroduces parameters and is forbidden.
\end{remark}

% ---------------------------------------------------------------------
\section{Route III: Graded Balanced-Ternary Map}
\label{sec:route3}

\begin{definition}[Graded Ternary Data]
\label{def:graded}
A \emph{grading} of the balanced-ternary map is a family of cocycles
\[
c_j\colon\mathbb{N}\to\{-1,0,+1\},\qquad j\in\{0,1,2\},
\]
together with the common underlying map \(T_3\), such that
\begin{equation}
\label{eq:graded-step}
T_3^{(j)}(n)\;:=\;T_3(n)
\quad\text{(set-map unchanged)},
\end{equation}
while observables use weights \(c_j\). The cocycles satisfy
\begin{equation}
\label{eq:cocycle-balance}
\sum_{j=0}^{2} c_j(n)\;=\;Q_{\mathrm{res}}\;=\;+1
\quad\text{for all \(n\) in the master tower}.
\end{equation}
\end{definition}

\begin{definition}[Contraction Rates and Resonant Periods]
\label{def:rates}
Write \(\lambda=\ln 3/539\) for the Banach contraction rate of the ungraded HQCC
map. Graded contractions \(\lambda_j\) must be rational multiples of \(\lambda\).
Graded resonant sub-periods \(\tau_j^{\mathrm{sub}}\) must be drawn from the
sub-harmonic set \(\{5,10,15,30,45\}\,\mathrm{s}\) (or integer commensurators of
\(G_4\))---never three independent fundamental periods.
\end{definition}

\begin{lemma}[Common Orbit Constraint]
\label{lem:common-orbit}
If the three generations are assigned three independent fundamental periods,
fixed-point rigidity (Lemma~\ref{lem:rigidity}) fails. Therefore any admissible
grading keeps a common termination depth \(\sigma=539\pm 1\) and a common
fundamental mode \(G_4\), with individuality in cocycle weights, residue classes
modulo 3, or sub-period structure.
\end{lemma}

\begin{theorem}[Route III Individuality]
\label{thm:route3}
Let \((c_0,c_1,c_2)\) be the canonical graded cocycles of
Section~\ref{sec:closed}. Then
\[
\iota_{III}\colon g_j\mapsto c_j
\]
is an admissible individuality structure: one orbit, one \(G_4\),
generation-dependent dynamical observables.
\end{theorem}

\begin{proof}[Proof sketch]
The set-map \(T_3\) and stopping time \(\sigma\) are generation-independent.
Inequivalence of the \(c_j\) (distinct mod-3 supports) breaks \(\mathrm{Sym}_3\).
Observables (mean sojourn near punctures, weighted transfer spectrum) depend on
\(j\) only through \(c_j\). All numerical inputs are chain integers.
\end{proof}

\begin{remark}[Failure Mode]
Independent fundamental periods not commensurate with \(G_4\) are inadmissible.
\end{remark}

% ---------------------------------------------------------------------
\section{Closed Integer Formulae}
\label{sec:closed}

All ``forced by the chain'' steps are fixed here by explicit integers from
\(\{3,7,11,13,61,539,1001,4880\}\). No continuous choice remains.

\begin{definition}[Canonical Step and Base Winding]
\label{def:Delta}
\[
\Delta \;:=\; |P| \;=\; 61,
\qquad
w_0 \;:=\; \sigma_{\mathrm{HQCC}} \;=\; 539.
\]
\end{definition}

\begin{lemma}[Canonical Windings]
\label{lem:canonical-w}
The unique arithmetic-progression winding assignment with base \(w_0=539\) and
step \(\Delta=61\) is
\begin{equation}
\label{eq:w-explicit}
w_j \;=\; 539 + 61\,j,\qquad j\in\{0,1,2\},
\end{equation}
i.e.
\[
(w_0,w_1,w_2) \;=\; (539,\,600,\,661).
\]
The total class is
\begin{equation}
\label{eq:Wtot}
W_{\mathrm{tot}} \;=\; 539+600+661 \;=\; 1800 \;=\; 3\times 600 \;=\; 3(w_0+\Delta).
\end{equation}
\textbf{Status (No-Go chapter).}
The pair \((w_0,\Delta)=(539,61)\) is \emph{conditional} on a non-circular supply
of \(\sigma=539\) (and of \(|P|=61\)). Flux democracy alone does not select
these anchors (Theorem on flux democracy no-go). By contrast, cocycles
\(c_j\) and shells \(k_j\in\{7,11,13\}\) remain a priori.
\end{lemma}

\begin{proof}
Substitute \eqref{eq:w-explicit} into \eqref{eq:AP-sum}. Summation is elementary.
Both \(539\) and \(61\) are Global Lemma Core integers; no other pair is used.
\end{proof}

\begin{lemma}[Canonical Shells]
\label{lem:canonical-k}
The unique unordered triple of prime shells forced by \(1001=7\times 11\times 13\)
is \(\{7,11,13\}\). Ordering by increasing prime (aligned with increasing
\(j\) and chirality of \(T\)) gives
\begin{equation}
\label{eq:k-explicit}
(k_0,k_1,k_2) \;=\; (7,\,11,\,13).
\end{equation}
\end{lemma}

\begin{definition}[Canonical Dictionaries \(\Phi,\Psi\)]
\label{def:PhiPsi}
\begin{align}
\Phi(w_j) &\;:=\; k_j \;=\; p_{j+1},
  \label{eq:Phi}\\
\Psi(k_j) &\;:=\; c_j,
  \label{eq:Psi}
\end{align}
where \(p_1=7\), \(p_2=11\), \(p_3=13\) are the prime factors of \(1001\) in
increasing order, and \(c_j\) is Definition~\ref{def:c-explicit}. Equivalently,
\[
j \;=\; \frac{w_j-539}{61} \;=\; \mathrm{index}(k_j\text{ in }(7,11,13)),
\]
so \(\Phi\) and \(\Psi\) are determined by the single generation index \(j\).
\end{definition}

\begin{definition}[Canonical Cocycles]
\label{def:c-explicit}
\begin{equation}
\label{eq:c-explicit}
c_j(n) \;:=\;
\begin{cases}
+1 & \text{if } n\equiv j \pmod{3},\\
0 & \text{otherwise}.
\end{cases}
\end{equation}
\end{definition}

\begin{lemma}[Cocycle Balance]
\label{lem:cocycle-ok}
For every \(n\in\mathbb{N}\),
\[
\sum_{j=0}^{2} c_j(n) \;=\; 1 \;=\; Q_{\mathrm{res}}.
\]
Thus \eqref{eq:cocycle-balance} holds pointwise on the master tower.
\end{lemma}

\begin{proof}
Exactly one residue class modulo 3 equals \(j\) for a unique \(j\in\{0,1,2\}\).
\end{proof}

\begin{definition}[Canonical Sub-Harmonics and Contractions]
\label{def:sub-explicit}
\begin{equation}
\label{eq:sub-explicit}
\bigl(\tau_0^{\mathrm{sub}},\tau_1^{\mathrm{sub}},\tau_2^{\mathrm{sub}}\bigr)
\;=\;
(5,\,15,\,45)\,\mathrm{s}
\;\subset\;
\{5,10,15,30,45\}\,\mathrm{s},
\end{equation}
ordered lightest \(\leftrightarrow\) shortest, and
\begin{equation}
\label{eq:lambda-explicit}
\lambda_j \;=\; \lambda \;=\; \frac{\ln 3}{539}
\quad\text{for all }j
\end{equation}
(common contraction; generation dependence lives in \(c_j\) and
\(\tau_j^{\mathrm{sub}}\) only).
\end{definition}

\begin{definition}[Canonical Lifetime Weight]
\label{def:f}
\begin{equation}
\label{eq:f}
f(k) \;:=\; \frac{1001}{k},
\qquad
\bigl(f(7),f(11),f(13)\bigr) \;=\; (143,\,91,\,77).
\end{equation}
\end{definition}

\begin{theorem}[Closure of All Forced Steps]
\label{thm:closure}
With Definitions~\ref{def:Delta}--\ref{def:f} and
Lemmas~\ref{lem:canonical-w}--\ref{lem:cocycle-ok}, every open ``forced by the
chain'' step in Routes I--III is an explicit integer formula. The only free
discrete choice retained is the global orientation
(increasing vs.\ decreasing \(j\)), fixed by the chirality of \(T\) to the
increasing convention \(j=0,1,2\leftrightarrow e,\mu,\tau\).
\end{theorem}

% ---------------------------------------------------------------------
\section{Composite Package}
\label{sec:composite}

\begin{definition}[Composite Individuality Structure]
\label{def:composite}
The \emph{composite package} is
\[
\iota \;=\; (\iota_I,\iota_{II},\iota_{III}),
\]
with
\begin{enumerate}[label=(\roman*)]
  \item Route I: \(g_j\mapsto w_j=539+61j\);
  \item Route II: \(w_j\mapsto k_j\in\{7,11,13\}\) via \(\Phi\);
  \item Route III: \(k_j\mapsto c_j\) via \(\Psi\), on the common \(T_3\) orbit.
\end{enumerate}
Consistency is the commuting square
\begin{equation}
\label{eq:dictionary}
\begin{array}{ccc}
\{g_j\} & \xrightarrow{\;\iota_I\;} & \{w_j\} \\[4pt]
\big\downarrow{\scriptstyle \iota_{III}} & & \big\downarrow{\scriptstyle \Phi} \\[4pt]
\{c_j\} & \xleftarrow{\;\Psi\;} & \{k_j\}
\end{array}
\end{equation}
with \(\Phi,\Psi\) as in Definition~\ref{def:PhiPsi}.
\end{definition}

\begin{theorem}[Individuality without a Second Axiom]
\label{thm:main}
The composite package \(\iota\) is an admissible individuality structure:
\begin{enumerate}[label=(\roman*)]
  \item Axiom 0 remains the sole primitive (cardinality three);
  \item the three generations are pairwise inequivalent as labeled individuals;
  \item \(G_4=539.9\,\mathrm{s}\), HQCC depth \(539\pm 1\), and \(|P|=61\) are
        unchanged;
  \item no free continuous parameter is introduced.
\end{enumerate}
\end{theorem}

\begin{proof}[Proof sketch]
Theorem~\ref{thm:closure} supplies explicit integer labels. Routes I--III are
admissible by Theorems~\ref{thm:route1}--\ref{thm:route3}. The dictionaries
\(\Phi,\Psi\) are bijections on a three-element set, hence introduce no continuum
of moduli. Conditions (i)--(iv) follow.
\end{proof}

\begin{corollary}[Symmetry Breaking Pattern]
\[
\mathrm{Sym}(\{0,1,2\})\;\longrightarrow\;\{1\},
\]
broken completely by \(\iota\). Partial activation of a single route may leave
a \(\mathbb{Z}_2\) stabilizer.
\end{corollary}

% ---------------------------------------------------------------------
\section{Canonical Generation-Label Table (\(e/\mu/\tau\))}
\label{sec:table}

\begin{definition}[Canonical Charged-Lepton Dictionary]
\label{def:leptons}
Identify
\[
(g_0,g_1,g_2) \;=\; (e,\,\mu,\,\tau)
\]
with the closed labels of Section~\ref{sec:closed}. The same \(j\) labels quark
generations when lifted: \((d,s,b)\) and \((u,c,t)\).
\end{definition}

\begin{table}[ht]
\centering
\textbf{Table: Canonical individuality labels for charged leptons}
\label{tab:labels}\\[0.5em]
\begin{tabular}{@{}llrrrrl@{}}
\toprule
Gen. & \(j\) & \(w_j\) & \(k_j\) & \(f(k_j)\) & \(\tau_j^{\mathrm{sub}}\) (s) & \(c_j\) support \\
\midrule
\(e\)   & 0 & 539 & 7  & 143 & 5  & \(n\equiv 0\pmod{3}\) \\
\(\mu\) & 1 & 600 & 11 & 91  & 15 & \(n\equiv 1\pmod{3}\) \\
\(\tau\)& 2 & 661 & 13 & 77  & 45 & \(n\equiv 2\pmod{3}\) \\
\bottomrule
\end{tabular}
\end{table}

\begin{table}[ht]
\centering
\textbf{Table: Global invariants preserved by the dictionary}
\label{tab:invariants}\\[0.5em]
\begin{tabular}{@{}ll@{}}
\toprule
Invariant & Value \\
\midrule
\(W_{\mathrm{tot}}=\sum_j w_j\) & \(1800=3\times 600\) \\
\(\prod_j k_j\) & \(7\times 11\times 13=1001\) \\
\(\sum_j c_j(n)\) & \(1=Q_{\mathrm{res}}=61\bmod 3\) \\
\(\lambda_j\) & \(\ln 3/539\) (common) \\
HQCC depth & \(539\pm 1\) \\
\(G_4\) & \(539.9\,\mathrm{s}\) \\
\(|P|\) & \(61\) \\
\bottomrule
\end{tabular}
\end{table}

\begin{corollary}[Explicit Lifetime Refinement]
Under \eqref{eq:tau-g} and \eqref{eq:f},
\begin{equation}
\label{eq:tau-explicit}
\tau_j \;\propto\; \frac{61}{1001}\cdot\frac{1001}{k_j}\cdot\frac{\phi^7}{539.9}
\;=\; \frac{61}{k_j}\cdot\frac{\phi^7}{539.9},
\end{equation}
so the shell weight cancels one factor of \(1001\) and leaves a pure prime in the
denominator: \(\tau_e\propto 61/7\), \(\tau_\mu\propto 61/11\),
\(\tau_\tau\propto 61/13\), times the common \(\phi^7/539.9\) factor. The standard
muon formula in the theorem chain corresponds to the \(j=1\) line with an overall
normalization \(\tau^-\) and resonance correction \(\delta_{\mathrm{res}}\).
\end{corollary}

\begin{corollary}[Discrete Mass Ordering]
Along Route I, winding gaps are constant:
\[
w_1-w_0 \;=\; w_2-w_1 \;=\; 61,
\]
so any observable monotone in \(w_j\) automatically orders
\(e\prec\mu\prec\tau\) (or the reverse if orientation is flipped---excluded by
chirality of \(T\)).
\end{corollary}

\begin{remark}[Quark Lift]
The same table applies to quarks by reinterpreting \(g_j\) as the \(j\)-th quark
generation. Flavour mixing is then the overlap of cocycles \(c_i,c_j\) on the
common HQCC orbit (CKM / PMNS as discrete cocycle correlators), with no
continuous Yukawa matrix beyond the finite dictionary.
\end{remark}

% ---------------------------------------------------------------------
\section{Comparison and Selection}
\label{sec:compare}

\begin{center}
\begin{tabular}{@{}llll@{}}
\toprule
Route & Carrier & Principal risk & Cleanest fit \\
\midrule
I & Windings on 11D cycles & Free winding family & Brane / flux language \\
II & Shells of \(\Delta^{1001}\) & Continuous radius & \(61/1001\) numerics \\
III & Graded \(T_3\) cocycles & Multiple fundamentals & HQCC core \\
Composite & \((w_j,k_j,c_j)\) & Over-constraining \(\Phi,\Psi\) & Full tower \\
\bottomrule
\end{tabular}
\end{center}

\begin{lemma}[Selection under Rigidity]
\label{lem:selection}
If a single route is preferred, Route III with common orbit and graded cocycle
is the least likely to reopen the fixed point of \(G_4\). If full geometric
language is required, the composite package is preferred, with Route III as the
dynamical readout and Routes I--II as geometric realization. The Axiomatic Book
takes the composite package with Table~\ref{tab:labels} as primary.
\end{lemma}

% ---------------------------------------------------------------------
\section{Falsifiability and Downstream Couplings}
\label{sec:falsify}

\begin{definition}[Individuality Observables]
Generation-dependent observables induced by \(\iota\) include:
\begin{enumerate}[label=(\roman*)]
  \item mass ratios \(m_i/m_j\) as functions of \((w_i-w_j)=61(i-j)\);
  \item lifetime ratios from \eqref{eq:tau-explicit}:
        \(\tau_i/\tau_j = k_j/k_i\) at fixed overall normalization;
  \item mixing angles as overlap of cocycles \(c_i,c_j\) on the common orbit;
  \item sub-harmonic phase offsets at \(\tau_j^{\mathrm{sub}}\in\{5,15,45\}\,\mathrm{s}\).
\end{enumerate}
\end{definition}

\begin{lemma}[Falsifiability]
\label{lem:falsify}
The canonical dictionary is falsifiable if measured mass or lifetime ratios
cannot be matched by Table~\ref{tab:labels} (or its single orientation reverse,
already excluded by chirality) while preserving \(G_4\) and \(|P|=61\). There is
no continuous retuning: the label set is finite.
\end{lemma}

\begin{corollary}[Lifetimeage Ratio Test]
At fixed normalization, \eqref{eq:tau-explicit} predicts
\[
\frac{\tau_\mu}{\tau_e} \;=\; \frac{7}{11},\qquad
\frac{\tau_\tau}{\tau_\mu} \;=\; \frac{11}{13},\qquad
\frac{\tau_\tau}{\tau_e} \;=\; \frac{7}{13},
\]
up to the known overall factors already present in the muon theorem line and
channel-dependent phase space. Gross violation of these discrete ratios (after
those factors) falsifies the canonical shell assignment.
\end{corollary}

% ---------------------------------------------------------------------
\section{Global Lemma Core Form}
\label{sec:glc}

\begin{lemma}[Individuality without a Second Axiom]
\label{lem:glc-individuality}
There exists a discrete, Axiom-0-forced assignment
\[
j\;\mapsto\; (w_j,k_j,c_j)
\;=\;
\bigl(539+61j,\; p_{j+1},\; \mathbf{1}_{n\equiv j\pmod{3}}\bigr)
\]
with \((p_1,p_2,p_3)=(7,11,13)\) such that the three generations are inequivalent
while \(G_4\), HQCC depth, and \(|P|\) remain unique.
\end{lemma}

% ---------------------------------------------------------------------
\section{Summary}
\label{sec:summary}

Axiom 0 yields three as a cardinality. Individuality is restored by the
composite integer dictionary:
\begin{enumerate}
  \item windings \(w_j=539+61j\) on 11-dimensional cycles;
  \item leftover \(+1\) on shells \(k_j\in\{7,11,13\}\) of \(\Delta^{1001}\);
  \item graded cocycles \(c_j=\mathbf{1}_{n\equiv j\pmod{3}}\) on a common \(T_3\)
        orbit of length \(539\pm 1\), with sub-harmonics \((5,15,45)\,\mathrm{s}\).
\end{enumerate}
For charged leptons, \(j=0,1,2\leftrightarrow e,\mu,\tau\)
(Table~\ref{tab:labels}). No second continuous axiom is required; fixed-point
rigidity is preserved; falsifiability is finite and discrete.

% end of chapter
  (after Axiom 0 / Math 1)
% =====================================================================

\chapter{Restoring Generation Individuality from Axiom 0}
\label{ch:individuality}

\section{Motivation and Problem Statement}

\begin{definition}[Multiplicity vs.\ Individuality]
\emph{Multiplicity} is the cardinality of the generation set. \emph{Individuality}
is a complete set of inequivalent labels distinguishing the three Standard-Model
generations as ordered individuals \((g_0,g_1,g_2)\), together with observables
(masses, mixings, lifetimes) that depend nontrivially on those labels.
\end{definition}

\begin{lemma}[Multiplicity from Axiom 0]
\label{lem:multiplicity}
Axiom 0 asserts that the universe contains exactly three generations of
Standard-Model fermions. This forces the qutrit structure, the non-perturbative
superpotential \(W_{np}=e^3\), the flux budget
\[
N_{\mathrm{flux}}=\lfloor e^3\times 3^5\rfloor=4880,
\]
HQCC termination in \(539\pm 1\) steps, the immutable breathing mode
\(G_4=539.9\,\mathrm{s}\), and the puncture count \(|P|=61\).
\end{lemma}

\begin{lemma}[Individuality Gap]
\label{lem:gap}
Axiom 0 alone determines multiplicity but does not determine individuality.
Any two bijections of the generation set are observationally equivalent at the
level of Axiom 0: the axiom is invariant under \(\mathrm{Sym}(\{0,1,2\})\).
\end{lemma}

\begin{corollary}[Need for Additional Discrete Structure]
To restore individuality while still beginning from the single empirical fact of
three generations, additional structure is required. The structure must be
\emph{discrete}, forced by integers already present in the derived chain
\(\{3,61,539,1001,4880,\ldots\}\), and must not reopen the fixed-point rigidity
of \(G_4\) or the puncture count \(|P|=61\).
\end{corollary}

\begin{definition}[Admissible Individuality Structure]
An \emph{admissible individuality structure} is a map
\[
\iota\colon \{g_0,g_1,g_2\}\longrightarrow \mathcal{L}
\]
into a discrete label space \(\mathcal{L}\) such that:
\begin{enumerate}[label=(\roman*)]
  \item \(\iota\) is injective (labels distinguish generations);
  \item \(\mathcal{L}\) is constructed only from Axiom 0 and the Global Lemma Core
        (no free continuous parameters);
  \item the push-forward of \(\iota\) leaves \(G_4=539.9\,\mathrm{s}\) and
        \(|P|=61\) invariant;
  \item residual \(\mathrm{Sym}(\{0,1,2\})\) is broken only by integers already
        fixed in the chain.
\end{enumerate}
\end{definition}

Three admissible routes are developed below and closed by explicit integer
formulae (Section~\ref{sec:closed}). They compose into a single package
(Section~\ref{sec:composite}) with a canonical charged-lepton table
(Section~\ref{sec:table}).

% ---------------------------------------------------------------------
\section{Shared Scaffolding}
\label{sec:scaffold}

\begin{definition}[Balanced-Ternary Syracuse Map]
\label{def:T3}
The balanced-ternary Syracuse map \(T_3\colon\mathbb{N}\to\mathbb{N}\) is the
integer-valued Collatz-type operator of the HQCC theorem (exact \(539\)-step
termination above the model threshold). Write \(\sigma(n)\) for the stopping
time to the fixed point \(1\).
\end{definition}

\begin{definition}[Residual \(+1\) Charge]
\label{def:residual}
The residual charge is the puncture class modulo the generation qutrit:
\begin{equation}
\label{eq:Qres}
Q_{\mathrm{res}}\;:=\; |P|\bmod 3 \;=\; 61\bmod 3 \;=\; 1.
\end{equation}
Equivalently, \(Q_{\mathrm{res}}=+1\in\mathbb{Z}/3\mathbb{Z}\) is the leftover unit
after balanced-ternary cancellation on the bulk cycles. (Note:
\(N_{\mathrm{flux}}=4880\equiv 2\pmod{3}\); the residual used for individuality is
the puncture residual \eqref{eq:Qres}, not the raw flux residue.)
\end{definition}

\begin{definition}[11-Dimensional Bulk and Cycles]
\label{def:11D}
Let \(X^{11}\) be the 11-dimensional bulk of the model, with temporal torsion
operator
\[
T=d+(\zeta^t-1)\wedge H_3,\qquad
\zeta=\exp\bigl(2\pi i/539.9\bigr).
\]
Let \(\mathcal{C}_k(X^{11})\) denote integral \(k\)-cycles relevant to flux
wrapping (in particular those supporting \(G_4\) and the 61 punctures on the
\(G_2\)-holonomy 7-manifold).
\end{definition}

\begin{definition}[1001-Simplex]
\label{def:1001}
Let \(\Delta^{1001}\) be the standard 1001-simplex. The integer
\[
1001=7\times 11\times 13
\]
is already present in the derived numerical chain (e.g.\ the factor \(61/1001\)
in lifetime formulae). Shells of \(\Delta^{1001}\) are discrete orbits under a
specified group action or filtration (Definition~\ref{def:shells}); they are
\emph{not} continuous radial coordinates.
\end{definition}

\begin{lemma}[Fixed-Point Rigidity Constraint]
\label{lem:rigidity}
Any admissible individuality structure must preserve the unique fixed point
\[
G_4=539.9\,\mathrm{s}
\]
satisfying simultaneously HQCC termination, flux quantization, puncture count,
entropy saturation, and the sub-harmonic set \(\{5,10,15,30,45\}\,\mathrm{s}\).
In particular, one may not introduce three independent fundamental periods.
\end{lemma}

\begin{definition}[Generation Index]
\label{def:index}
Write \(j\in\{0,1,2\}\) for the generation index, with the charged-lepton
identification of Section~\ref{sec:table}:
\[
j=0\leftrightarrow e,\qquad
j=1\leftrightarrow \mu,\qquad
j=2\leftrightarrow \tau.
\]
The same index labels the three quark generations \((d,s,b)\) and \((u,c,t)\)
when individuality is lifted from leptons to quarks.
\end{definition}

% ---------------------------------------------------------------------
\section{Route I: Topological Charges and Winding Numbers}
\label{sec:route1}

\begin{definition}[Generation Winding Data]
\label{def:windings}
A \emph{generation winding assignment} is a triple
\[
\mathbf{w}=(w_0,w_1,w_2)\in H_k\bigl(X^{11};\mathbb{Z}\bigr)^3
\]
(or integer winding numbers of three distinguished cycles
\(\gamma_0,\gamma_1,\gamma_2\in\mathcal{C}_k(X^{11})\)) such that
\begin{equation}
\label{eq:winding-sum}
\sum_{j=0}^{2} w_j \;=\; W_{\mathrm{tot}},
\end{equation}
where \(W_{\mathrm{tot}}\) is an integer fixed by the closed formulae of
Section~\ref{sec:closed} (no free continuous modulus).
\end{definition}

\begin{definition}[Arithmetic Progression Form]
\label{def:AP}
The assignment \(\mathbf{w}\) is in \emph{arithmetic progression form} if there
exist integers \(w_0,\Delta\) with
\begin{equation}
\label{eq:AP}
w_j \;=\; w_0 + j\,\Delta,\qquad j\in\{0,1,2\},
\end{equation}
and hence
\begin{equation}
\label{eq:AP-sum}
W_{\mathrm{tot}} \;=\; 3w_0 + 3\Delta \;=\; 3(w_0+\Delta).
\end{equation}
\end{definition}

\begin{theorem}[Route I Individuality]
\label{thm:route1}
Suppose \(\mathbf{w}\) is the canonical arithmetic-progression assignment of
Section~\ref{sec:closed}. Then
\[
\iota_I\colon g_j\mapsto w_j
\]
is an admissible individuality structure. Generations are distinguished by
inequivalent topological charges on the 11-dimensional cycles, while
multiplicity and global flux remain Axiom-0 data.
\end{theorem}

\begin{proof}[Proof sketch]
Injectivity follows from \(\Delta=61\neq 0\). Discreteness is immediate from
integral homology. Invariance of \(G_4\) and \(|P|\) follows because
\eqref{eq:winding-sum} preserves the total winding class used for the breathing
mode and puncture count; only the distribution among three cycles changes.
Residual \(\mathrm{Sym}_3\) is broken by the ordered progression
\(w_0<w_1<w_2\), with orientation fixed by the chirality of \(T\).
\end{proof}

\begin{corollary}[Mass Hierarchy from Wrapping]
If particle masses scale with wrapping volume (or with the period of \(T\) along
\(\gamma_j\)), then
\[
m_j \;\propto\; \mathrm{Vol}(\gamma_j)
\quad\text{or}\quad
m_j \;\propto\; \bigl|\langle T,[ \gamma_j]\rangle\bigr|
\]
yields a discrete hierarchy indexed by \(w_j\), with no free continuous Yukawa
modulus beyond integers already in the chain.
\end{corollary}

\begin{remark}[Failure Mode]
If \(\mathbf{w}\) is chosen from an open family of windings not constrained by
\eqref{eq:winding-sum} and Section~\ref{sec:closed}, individuality is smuggled
in as a second axiom. Route I is admissible only in the forced form.
\end{remark}

% ---------------------------------------------------------------------
\section{Route II: Leftover \(+1\) on Shells of the 1001-Simplex}
\label{sec:route2}

\begin{definition}[Shells of \(\Delta^{1001}\)]
\label{def:shells}
Let \(\Gamma=\mathbb{Z}_7\times\mathbb{Z}_{11}\times\mathbb{Z}_{13}\) act by the
prime factorization \(1001=7\times 11\times 13\) (equivalently,
\(\Gamma_3=\mathbb{Z}_3\) by the generation qutrit). A \emph{shell} is a
\(\Gamma\)-orbit or a grade of a \(\Gamma\)-equivariant filtration
\[
\Delta^{1001}\;=\;\bigsqcup_{\ell} S_\ell.
\]
Shell indices are integers; there is no continuous radial coordinate.
\end{definition}

\begin{definition}[Generation-Dependent Embedding of \(Q_{\mathrm{res}}\)]
\label{def:embed}
A \emph{residual embedding} assigns the leftover charge \(Q_{\mathrm{res}}=+1\)
to three inequivalent shells:
\begin{equation}
\label{eq:shells}
\iota_{II}\colon g_j \mapsto S_{k_j},\qquad
k_0,k_1,k_2\text{ pairwise distinct}.
\end{equation}
\end{definition}

\begin{lemma}[Single-Shell Obstruction]
\label{lem:single-shell}
If \(Q_{\mathrm{res}}=+1\) is placed in a single shell only, then either
\begin{enumerate}[label=(\alph*)]
  \item the three-generation qutrit collapses to a single effective grade
        (loss of multiplicity), or
  \item flux quantization / entropy conservation across the \(+U/-U\) interface
        fails to match the Global Lemma Core.
\end{enumerate}
Hence any residual embedding compatible with Axiom 0 must occupy (at least)
three inequivalent shells.
\end{lemma}

\begin{theorem}[Route II Individuality]
\label{thm:route2}
Let \(\iota_{II}\) be the canonical residual embedding of
Section~\ref{sec:closed} with shell indices \(\{7,11,13\}\). Then \(\iota_{II}\)
is an admissible individuality structure.
\end{theorem}

\begin{proof}[Proof sketch]
Injectivity is pairwise distinct primes. Discreteness is by
Definition~\ref{def:shells}. The total residual remains \(+1\), so \(G_4\) and
\(|P|\) are unchanged. The identity \(7\times 11\times 13=1001\) and the factor
\(61/1001\) already appear in lifetime formulae.
\end{proof}

\begin{corollary}[Lifetimeage Coupling]
The model lifetime schematic
\[
\tau \;\propto\; \frac{61}{1001}\cdot\frac{\phi^7}{539.9}
\]
admits the generation refinement
\begin{equation}
\label{eq:tau-g}
\tau_j \;\propto\; \frac{61}{1001}\cdot f(k_j)\cdot\frac{\phi^7}{539.9},
\end{equation}
with the fixed rational weight of Section~\ref{sec:closed}
(no free continuous coupling).
\end{corollary}

\begin{remark}[Failure Mode]
Continuous ``shell radius'' reintroduces parameters and is forbidden.
\end{remark}

% ---------------------------------------------------------------------
\section{Route III: Graded Balanced-Ternary Map}
\label{sec:route3}

\begin{definition}[Graded Ternary Data]
\label{def:graded}
A \emph{grading} of the balanced-ternary map is a family of cocycles
\[
c_j\colon\mathbb{N}\to\{-1,0,+1\},\qquad j\in\{0,1,2\},
\]
together with the common underlying map \(T_3\), such that
\begin{equation}
\label{eq:graded-step}
T_3^{(j)}(n)\;:=\;T_3(n)
\quad\text{(set-map unchanged)},
\end{equation}
while observables use weights \(c_j\). The cocycles satisfy
\begin{equation}
\label{eq:cocycle-balance}
\sum_{j=0}^{2} c_j(n)\;=\;Q_{\mathrm{res}}\;=\;+1
\quad\text{for all \(n\) in the master tower}.
\end{equation}
\end{definition}

\begin{definition}[Contraction Rates and Resonant Periods]
\label{def:rates}
Write \(\lambda=\ln 3/539\) for the Banach contraction rate of the ungraded HQCC
map. Graded contractions \(\lambda_j\) must be rational multiples of \(\lambda\).
Graded resonant sub-periods \(\tau_j^{\mathrm{sub}}\) must be drawn from the
sub-harmonic set \(\{5,10,15,30,45\}\,\mathrm{s}\) (or integer commensurators of
\(G_4\))---never three independent fundamental periods.
\end{definition}

\begin{lemma}[Common Orbit Constraint]
\label{lem:common-orbit}
If the three generations are assigned three independent fundamental periods,
fixed-point rigidity (Lemma~\ref{lem:rigidity}) fails. Therefore any admissible
grading keeps a common termination depth \(\sigma=539\pm 1\) and a common
fundamental mode \(G_4\), with individuality in cocycle weights, residue classes
modulo 3, or sub-period structure.
\end{lemma}

\begin{theorem}[Route III Individuality]
\label{thm:route3}
Let \((c_0,c_1,c_2)\) be the canonical graded cocycles of
Section~\ref{sec:closed}. Then
\[
\iota_{III}\colon g_j\mapsto c_j
\]
is an admissible individuality structure: one orbit, one \(G_4\),
generation-dependent dynamical observables.
\end{theorem}

\begin{proof}[Proof sketch]
The set-map \(T_3\) and stopping time \(\sigma\) are generation-independent.
Inequivalence of the \(c_j\) (distinct mod-3 supports) breaks \(\mathrm{Sym}_3\).
Observables (mean sojourn near punctures, weighted transfer spectrum) depend on
\(j\) only through \(c_j\). All numerical inputs are chain integers.
\end{proof}

\begin{remark}[Failure Mode]
Independent fundamental periods not commensurate with \(G_4\) are inadmissible.
\end{remark}

% ---------------------------------------------------------------------
\section{Closed Integer Formulae}
\label{sec:closed}

All ``forced by the chain'' steps are fixed here by explicit integers from
\(\{3,7,11,13,61,539,1001,4880\}\). No continuous choice remains.

\begin{definition}[Canonical Step and Base Winding]
\label{def:Delta}
\[
\Delta \;:=\; |P| \;=\; 61,
\qquad
w_0 \;:=\; \sigma_{\mathrm{HQCC}} \;=\; 539.
\]
\end{definition}

\begin{lemma}[Canonical Windings]
\label{lem:canonical-w}
The unique arithmetic-progression winding assignment with base \(w_0=539\) and
step \(\Delta=61\) is
\begin{equation}
\label{eq:w-explicit}
w_j \;=\; 539 + 61\,j,\qquad j\in\{0,1,2\},
\end{equation}
i.e.
\[
(w_0,w_1,w_2) \;=\; (539,\,600,\,661).
\]
The total class is
\begin{equation}
\label{eq:Wtot}
W_{\mathrm{tot}} \;=\; 539+600+661 \;=\; 1800 \;=\; 3\times 600 \;=\; 3(w_0+\Delta).
\end{equation}
\textbf{Status (No-Go chapter).}
The pair \((w_0,\Delta)=(539,61)\) is \emph{conditional} on a non-circular supply
of \(\sigma=539\) (and of \(|P|=61\)). Flux democracy alone does not select
these anchors (Theorem on flux democracy no-go). By contrast, cocycles
\(c_j\) and shells \(k_j\in\{7,11,13\}\) remain a priori.
\end{lemma}

\begin{proof}
Substitute \eqref{eq:w-explicit} into \eqref{eq:AP-sum}. Summation is elementary.
Both \(539\) and \(61\) are Global Lemma Core integers; no other pair is used.
\end{proof}

\begin{lemma}[Canonical Shells]
\label{lem:canonical-k}
The unique unordered triple of prime shells forced by \(1001=7\times 11\times 13\)
is \(\{7,11,13\}\). Ordering by increasing prime (aligned with increasing
\(j\) and chirality of \(T\)) gives
\begin{equation}
\label{eq:k-explicit}
(k_0,k_1,k_2) \;=\; (7,\,11,\,13).
\end{equation}
\end{lemma}

\begin{definition}[Canonical Dictionaries \(\Phi,\Psi\)]
\label{def:PhiPsi}
\begin{align}
\Phi(w_j) &\;:=\; k_j \;=\; p_{j+1},
  \label{eq:Phi}\\
\Psi(k_j) &\;:=\; c_j,
  \label{eq:Psi}
\end{align}
where \(p_1=7\), \(p_2=11\), \(p_3=13\) are the prime factors of \(1001\) in
increasing order, and \(c_j\) is Definition~\ref{def:c-explicit}. Equivalently,
\[
j \;=\; \frac{w_j-539}{61} \;=\; \mathrm{index}(k_j\text{ in }(7,11,13)),
\]
so \(\Phi\) and \(\Psi\) are determined by the single generation index \(j\).
\end{definition}

\begin{definition}[Canonical Cocycles]
\label{def:c-explicit}
\begin{equation}
\label{eq:c-explicit}
c_j(n) \;:=\;
\begin{cases}
+1 & \text{if } n\equiv j \pmod{3},\\
0 & \text{otherwise}.
\end{cases}
\end{equation}
\end{definition}

\begin{lemma}[Cocycle Balance]
\label{lem:cocycle-ok}
For every \(n\in\mathbb{N}\),
\[
\sum_{j=0}^{2} c_j(n) \;=\; 1 \;=\; Q_{\mathrm{res}}.
\]
Thus \eqref{eq:cocycle-balance} holds pointwise on the master tower.
\end{lemma}

\begin{proof}
Exactly one residue class modulo 3 equals \(j\) for a unique \(j\in\{0,1,2\}\).
\end{proof}

\begin{definition}[Canonical Sub-Harmonics and Contractions]
\label{def:sub-explicit}
\begin{equation}
\label{eq:sub-explicit}
\bigl(\tau_0^{\mathrm{sub}},\tau_1^{\mathrm{sub}},\tau_2^{\mathrm{sub}}\bigr)
\;=\;
(5,\,15,\,45)\,\mathrm{s}
\;\subset\;
\{5,10,15,30,45\}\,\mathrm{s},
\end{equation}
ordered lightest \(\leftrightarrow\) shortest, and
\begin{equation}
\label{eq:lambda-explicit}
\lambda_j \;=\; \lambda \;=\; \frac{\ln 3}{539}
\quad\text{for all }j
\end{equation}
(common contraction; generation dependence lives in \(c_j\) and
\(\tau_j^{\mathrm{sub}}\) only).
\end{definition}

\begin{definition}[Canonical Lifetime Weight]
\label{def:f}
\begin{equation}
\label{eq:f}
f(k) \;:=\; \frac{1001}{k},
\qquad
\bigl(f(7),f(11),f(13)\bigr) \;=\; (143,\,91,\,77).
\end{equation}
\end{definition}

\begin{theorem}[Closure of All Forced Steps]
\label{thm:closure}
With Definitions~\ref{def:Delta}--\ref{def:f} and
Lemmas~\ref{lem:canonical-w}--\ref{lem:cocycle-ok}, every open ``forced by the
chain'' step in Routes I--III is an explicit integer formula. The only free
discrete choice retained is the global orientation
(increasing vs.\ decreasing \(j\)), fixed by the chirality of \(T\) to the
increasing convention \(j=0,1,2\leftrightarrow e,\mu,\tau\).
\end{theorem}

% ---------------------------------------------------------------------
\section{Composite Package}
\label{sec:composite}

\begin{definition}[Composite Individuality Structure]
\label{def:composite}
The \emph{composite package} is
\[
\iota \;=\; (\iota_I,\iota_{II},\iota_{III}),
\]
with
\begin{enumerate}[label=(\roman*)]
  \item Route I: \(g_j\mapsto w_j=539+61j\);
  \item Route II: \(w_j\mapsto k_j\in\{7,11,13\}\) via \(\Phi\);
  \item Route III: \(k_j\mapsto c_j\) via \(\Psi\), on the common \(T_3\) orbit.
\end{enumerate}
Consistency is the commuting square
\begin{equation}
\label{eq:dictionary}
\begin{array}{ccc}
\{g_j\} & \xrightarrow{\;\iota_I\;} & \{w_j\} \\[4pt]
\big\downarrow{\scriptstyle \iota_{III}} & & \big\downarrow{\scriptstyle \Phi} \\[4pt]
\{c_j\} & \xleftarrow{\;\Psi\;} & \{k_j\}
\end{array}
\end{equation}
with \(\Phi,\Psi\) as in Definition~\ref{def:PhiPsi}.
\end{definition}

\begin{theorem}[Individuality without a Second Axiom]
\label{thm:main}
The composite package \(\iota\) is an admissible individuality structure:
\begin{enumerate}[label=(\roman*)]
  \item Axiom 0 remains the sole primitive (cardinality three);
  \item the three generations are pairwise inequivalent as labeled individuals;
  \item \(G_4=539.9\,\mathrm{s}\), HQCC depth \(539\pm 1\), and \(|P|=61\) are
        unchanged;
  \item no free continuous parameter is introduced.
\end{enumerate}
\end{theorem}

\begin{proof}[Proof sketch]
Theorem~\ref{thm:closure} supplies explicit integer labels. Routes I--III are
admissible by Theorems~\ref{thm:route1}--\ref{thm:route3}. The dictionaries
\(\Phi,\Psi\) are bijections on a three-element set, hence introduce no continuum
of moduli. Conditions (i)--(iv) follow.
\end{proof}

\begin{corollary}[Symmetry Breaking Pattern]
\[
\mathrm{Sym}(\{0,1,2\})\;\longrightarrow\;\{1\},
\]
broken completely by \(\iota\). Partial activation of a single route may leave
a \(\mathbb{Z}_2\) stabilizer.
\end{corollary}

% ---------------------------------------------------------------------
\section{Canonical Generation-Label Table (\(e/\mu/\tau\))}
\label{sec:table}

\begin{definition}[Canonical Charged-Lepton Dictionary]
\label{def:leptons}
Identify
\[
(g_0,g_1,g_2) \;=\; (e,\,\mu,\,\tau)
\]
with the closed labels of Section~\ref{sec:closed}. The same \(j\) labels quark
generations when lifted: \((d,s,b)\) and \((u,c,t)\).
\end{definition}

\begin{table}[ht]
\centering
\textbf{Table: Canonical individuality labels for charged leptons}
\label{tab:labels}\\[0.5em]
\begin{tabular}{@{}llrrrrl@{}}
\toprule
Gen. & \(j\) & \(w_j\) & \(k_j\) & \(f(k_j)\) & \(\tau_j^{\mathrm{sub}}\) (s) & \(c_j\) support \\
\midrule
\(e\)   & 0 & 539 & 7  & 143 & 5  & \(n\equiv 0\pmod{3}\) \\
\(\mu\) & 1 & 600 & 11 & 91  & 15 & \(n\equiv 1\pmod{3}\) \\
\(\tau\)& 2 & 661 & 13 & 77  & 45 & \(n\equiv 2\pmod{3}\) \\
\bottomrule
\end{tabular}
\end{table}

\begin{table}[ht]
\centering
\textbf{Table: Global invariants preserved by the dictionary}
\label{tab:invariants}\\[0.5em]
\begin{tabular}{@{}ll@{}}
\toprule
Invariant & Value \\
\midrule
\(W_{\mathrm{tot}}=\sum_j w_j\) & \(1800=3\times 600\) \\
\(\prod_j k_j\) & \(7\times 11\times 13=1001\) \\
\(\sum_j c_j(n)\) & \(1=Q_{\mathrm{res}}=61\bmod 3\) \\
\(\lambda_j\) & \(\ln 3/539\) (common) \\
HQCC depth & \(539\pm 1\) \\
\(G_4\) & \(539.9\,\mathrm{s}\) \\
\(|P|\) & \(61\) \\
\bottomrule
\end{tabular}
\end{table}

\begin{corollary}[Explicit Lifetime Refinement]
Under \eqref{eq:tau-g} and \eqref{eq:f},
\begin{equation}
\label{eq:tau-explicit}
\tau_j \;\propto\; \frac{61}{1001}\cdot\frac{1001}{k_j}\cdot\frac{\phi^7}{539.9}
\;=\; \frac{61}{k_j}\cdot\frac{\phi^7}{539.9},
\end{equation}
so the shell weight cancels one factor of \(1001\) and leaves a pure prime in the
denominator: \(\tau_e\propto 61/7\), \(\tau_\mu\propto 61/11\),
\(\tau_\tau\propto 61/13\), times the common \(\phi^7/539.9\) factor. The standard
muon formula in the theorem chain corresponds to the \(j=1\) line with an overall
normalization \(\tau^-\) and resonance correction \(\delta_{\mathrm{res}}\).
\end{corollary}

\begin{corollary}[Discrete Mass Ordering]
Along Route I, winding gaps are constant:
\[
w_1-w_0 \;=\; w_2-w_1 \;=\; 61,
\]
so any observable monotone in \(w_j\) automatically orders
\(e\prec\mu\prec\tau\) (or the reverse if orientation is flipped---excluded by
chirality of \(T\)).
\end{corollary}

\begin{remark}[Quark Lift]
The same table applies to quarks by reinterpreting \(g_j\) as the \(j\)-th quark
generation. Flavour mixing is then the overlap of cocycles \(c_i,c_j\) on the
common HQCC orbit (CKM / PMNS as discrete cocycle correlators), with no
continuous Yukawa matrix beyond the finite dictionary.
\end{remark}

% ---------------------------------------------------------------------
\section{Comparison and Selection}
\label{sec:compare}

\begin{center}
\begin{tabular}{@{}llll@{}}
\toprule
Route & Carrier & Principal risk & Cleanest fit \\
\midrule
I & Windings on 11D cycles & Free winding family & Brane / flux language \\
II & Shells of \(\Delta^{1001}\) & Continuous radius & \(61/1001\) numerics \\
III & Graded \(T_3\) cocycles & Multiple fundamentals & HQCC core \\
Composite & \((w_j,k_j,c_j)\) & Over-constraining \(\Phi,\Psi\) & Full tower \\
\bottomrule
\end{tabular}
\end{center}

\begin{lemma}[Selection under Rigidity]
\label{lem:selection}
If a single route is preferred, Route III with common orbit and graded cocycle
is the least likely to reopen the fixed point of \(G_4\). If full geometric
language is required, the composite package is preferred, with Route III as the
dynamical readout and Routes I--II as geometric realization. The Axiomatic Book
takes the composite package with Table~\ref{tab:labels} as primary.
\end{lemma}

% ---------------------------------------------------------------------
\section{Falsifiability and Downstream Couplings}
\label{sec:falsify}

\begin{definition}[Individuality Observables]
Generation-dependent observables induced by \(\iota\) include:
\begin{enumerate}[label=(\roman*)]
  \item mass ratios \(m_i/m_j\) as functions of \((w_i-w_j)=61(i-j)\);
  \item lifetime ratios from \eqref{eq:tau-explicit}:
        \(\tau_i/\tau_j = k_j/k_i\) at fixed overall normalization;
  \item mixing angles as overlap of cocycles \(c_i,c_j\) on the common orbit;
  \item sub-harmonic phase offsets at \(\tau_j^{\mathrm{sub}}\in\{5,15,45\}\,\mathrm{s}\).
\end{enumerate}
\end{definition}

\begin{lemma}[Falsifiability]
\label{lem:falsify}
The canonical dictionary is falsifiable if measured mass or lifetime ratios
cannot be matched by Table~\ref{tab:labels} (or its single orientation reverse,
already excluded by chirality) while preserving \(G_4\) and \(|P|=61\). There is
no continuous retuning: the label set is finite.
\end{lemma}

\begin{corollary}[Lifetimeage Ratio Test]
At fixed normalization, \eqref{eq:tau-explicit} predicts
\[
\frac{\tau_\mu}{\tau_e} \;=\; \frac{7}{11},\qquad
\frac{\tau_\tau}{\tau_\mu} \;=\; \frac{11}{13},\qquad
\frac{\tau_\tau}{\tau_e} \;=\; \frac{7}{13},
\]
up to the known overall factors already present in the muon theorem line and
channel-dependent phase space. Gross violation of these discrete ratios (after
those factors) falsifies the canonical shell assignment.
\end{corollary}

% ---------------------------------------------------------------------
\section{Global Lemma Core Form}
\label{sec:glc}

\begin{lemma}[Individuality without a Second Axiom]
\label{lem:glc-individuality}
There exists a discrete, Axiom-0-forced assignment
\[
j\;\mapsto\; (w_j,k_j,c_j)
\;=\;
\bigl(539+61j,\; p_{j+1},\; \mathbf{1}_{n\equiv j\pmod{3}}\bigr)
\]
with \((p_1,p_2,p_3)=(7,11,13)\) such that the three generations are inequivalent
while \(G_4\), HQCC depth, and \(|P|\) remain unique.
\end{lemma}

% ---------------------------------------------------------------------
\section{Summary}
\label{sec:summary}

Axiom 0 yields three as a cardinality. Individuality is restored by the
composite integer dictionary:
\begin{enumerate}
  \item windings \(w_j=539+61j\) on 11-dimensional cycles;
  \item leftover \(+1\) on shells \(k_j\in\{7,11,13\}\) of \(\Delta^{1001}\);
  \item graded cocycles \(c_j=\mathbf{1}_{n\equiv j\pmod{3}}\) on a common \(T_3\)
        orbit of length \(539\pm 1\), with sub-harmonics \((5,15,45)\,\mathrm{s}\).
\end{enumerate}
For charged leptons, \(j=0,1,2\leftrightarrow e,\mu,\tau\)
(Table~\ref{tab:labels}). No second continuous axiom is required; fixed-point
rigidity is preserved; falsifiability is finite and discrete.

% end of chapter
  (after Axiom 0 / Math 1)
% =====================================================================

\chapter{Restoring Generation Individuality from Axiom 0}
\label{ch:individuality}

\section{Motivation and Problem Statement}

\begin{definition}[Multiplicity vs.\ Individuality]
\emph{Multiplicity} is the cardinality of the generation set. \emph{Individuality}
is a complete set of inequivalent labels distinguishing the three Standard-Model
generations as ordered individuals \((g_0,g_1,g_2)\), together with observables
(masses, mixings, lifetimes) that depend nontrivially on those labels.
\end{definition}

\begin{lemma}[Multiplicity from Axiom 0]
\label{lem:multiplicity}
Axiom 0 asserts that the universe contains exactly three generations of
Standard-Model fermions. This forces the qutrit structure, the non-perturbative
superpotential \(W_{np}=e^3\), the flux budget
\[
N_{\mathrm{flux}}=\lfloor e^3\times 3^5\rfloor=4880,
\]
HQCC termination in \(539\pm 1\) steps, the immutable breathing mode
\(G_4=539.9\,\mathrm{s}\), and the puncture count \(|P|=61\).
\end{lemma}

\begin{lemma}[Individuality Gap]
\label{lem:gap}
Axiom 0 alone determines multiplicity but does not determine individuality.
Any two bijections of the generation set are observationally equivalent at the
level of Axiom 0: the axiom is invariant under \(\mathrm{Sym}(\{0,1,2\})\).
\end{lemma}

\begin{corollary}[Need for Additional Discrete Structure]
To restore individuality while still beginning from the single empirical fact of
three generations, additional structure is required. The structure must be
\emph{discrete}, forced by integers already present in the derived chain
\(\{3,61,539,1001,4880,\ldots\}\), and must not reopen the fixed-point rigidity
of \(G_4\) or the puncture count \(|P|=61\).
\end{corollary}

\begin{definition}[Admissible Individuality Structure]
An \emph{admissible individuality structure} is a map
\[
\iota\colon \{g_0,g_1,g_2\}\longrightarrow \mathcal{L}
\]
into a discrete label space \(\mathcal{L}\) such that:
\begin{enumerate}[label=(\roman*)]
  \item \(\iota\) is injective (labels distinguish generations);
  \item \(\mathcal{L}\) is constructed only from Axiom 0 and the Global Lemma Core
        (no free continuous parameters);
  \item the push-forward of \(\iota\) leaves \(G_4=539.9\,\mathrm{s}\) and
        \(|P|=61\) invariant;
  \item residual \(\mathrm{Sym}(\{0,1,2\})\) is broken only by integers already
        fixed in the chain.
\end{enumerate}
\end{definition}

Three admissible routes are developed below and closed by explicit integer
formulae (Section~\ref{sec:closed}). They compose into a single package
(Section~\ref{sec:composite}) with a canonical charged-lepton table
(Section~\ref{sec:table}).

% ---------------------------------------------------------------------
\section{Shared Scaffolding}
\label{sec:scaffold}

\begin{definition}[Balanced-Ternary Syracuse Map]
\label{def:T3}
The balanced-ternary Syracuse map \(T_3\colon\mathbb{N}\to\mathbb{N}\) is the
integer-valued Collatz-type operator of the HQCC theorem (exact \(539\)-step
termination above the model threshold). Write \(\sigma(n)\) for the stopping
time to the fixed point \(1\).
\end{definition}

\begin{definition}[Residual \(+1\) Charge]
\label{def:residual}
The residual charge is the puncture class modulo the generation qutrit:
\begin{equation}
\label{eq:Qres}
Q_{\mathrm{res}}\;:=\; |P|\bmod 3 \;=\; 61\bmod 3 \;=\; 1.
\end{equation}
Equivalently, \(Q_{\mathrm{res}}=+1\in\mathbb{Z}/3\mathbb{Z}\) is the leftover unit
after balanced-ternary cancellation on the bulk cycles. (Note:
\(N_{\mathrm{flux}}=4880\equiv 2\pmod{3}\); the residual used for individuality is
the puncture residual \eqref{eq:Qres}, not the raw flux residue.)
\end{definition}

\begin{definition}[11-Dimensional Bulk and Cycles]
\label{def:11D}
Let \(X^{11}\) be the 11-dimensional bulk of the model, with temporal torsion
operator
\[
T=d+(\zeta^t-1)\wedge H_3,\qquad
\zeta=\exp\bigl(2\pi i/539.9\bigr).
\]
Let \(\mathcal{C}_k(X^{11})\) denote integral \(k\)-cycles relevant to flux
wrapping (in particular those supporting \(G_4\) and the 61 punctures on the
\(G_2\)-holonomy 7-manifold).
\end{definition}

\begin{definition}[1001-Simplex]
\label{def:1001}
Let \(\Delta^{1001}\) be the standard 1001-simplex. The integer
\[
1001=7\times 11\times 13
\]
is already present in the derived numerical chain (e.g.\ the factor \(61/1001\)
in lifetime formulae). Shells of \(\Delta^{1001}\) are discrete orbits under a
specified group action or filtration (Definition~\ref{def:shells}); they are
\emph{not} continuous radial coordinates.
\end{definition}

\begin{lemma}[Fixed-Point Rigidity Constraint]
\label{lem:rigidity}
Any admissible individuality structure must preserve the unique fixed point
\[
G_4=539.9\,\mathrm{s}
\]
satisfying simultaneously HQCC termination, flux quantization, puncture count,
entropy saturation, and the sub-harmonic set \(\{5,10,15,30,45\}\,\mathrm{s}\).
In particular, one may not introduce three independent fundamental periods.
\end{lemma}

\begin{definition}[Generation Index]
\label{def:index}
Write \(j\in\{0,1,2\}\) for the generation index, with the charged-lepton
identification of Section~\ref{sec:table}:
\[
j=0\leftrightarrow e,\qquad
j=1\leftrightarrow \mu,\qquad
j=2\leftrightarrow \tau.
\]
The same index labels the three quark generations \((d,s,b)\) and \((u,c,t)\)
when individuality is lifted from leptons to quarks.
\end{definition}

% ---------------------------------------------------------------------
\section{Route I: Topological Charges and Winding Numbers}
\label{sec:route1}

\begin{definition}[Generation Winding Data]
\label{def:windings}
A \emph{generation winding assignment} is a triple
\[
\mathbf{w}=(w_0,w_1,w_2)\in H_k\bigl(X^{11};\mathbb{Z}\bigr)^3
\]
(or integer winding numbers of three distinguished cycles
\(\gamma_0,\gamma_1,\gamma_2\in\mathcal{C}_k(X^{11})\)) such that
\begin{equation}
\label{eq:winding-sum}
\sum_{j=0}^{2} w_j \;=\; W_{\mathrm{tot}},
\end{equation}
where \(W_{\mathrm{tot}}\) is an integer fixed by the closed formulae of
Section~\ref{sec:closed} (no free continuous modulus).
\end{definition}

\begin{definition}[Arithmetic Progression Form]
\label{def:AP}
The assignment \(\mathbf{w}\) is in \emph{arithmetic progression form} if there
exist integers \(w_0,\Delta\) with
\begin{equation}
\label{eq:AP}
w_j \;=\; w_0 + j\,\Delta,\qquad j\in\{0,1,2\},
\end{equation}
and hence
\begin{equation}
\label{eq:AP-sum}
W_{\mathrm{tot}} \;=\; 3w_0 + 3\Delta \;=\; 3(w_0+\Delta).
\end{equation}
\end{definition}

\begin{theorem}[Route I Individuality]
\label{thm:route1}
Suppose \(\mathbf{w}\) is the canonical arithmetic-progression assignment of
Section~\ref{sec:closed}. Then
\[
\iota_I\colon g_j\mapsto w_j
\]
is an admissible individuality structure. Generations are distinguished by
inequivalent topological charges on the 11-dimensional cycles, while
multiplicity and global flux remain Axiom-0 data.
\end{theorem}

\begin{proof}[Proof sketch]
Injectivity follows from \(\Delta=61\neq 0\). Discreteness is immediate from
integral homology. Invariance of \(G_4\) and \(|P|\) follows because
\eqref{eq:winding-sum} preserves the total winding class used for the breathing
mode and puncture count; only the distribution among three cycles changes.
Residual \(\mathrm{Sym}_3\) is broken by the ordered progression
\(w_0<w_1<w_2\), with orientation fixed by the chirality of \(T\).
\end{proof}

\begin{corollary}[Mass Hierarchy from Wrapping]
If particle masses scale with wrapping volume (or with the period of \(T\) along
\(\gamma_j\)), then
\[
m_j \;\propto\; \mathrm{Vol}(\gamma_j)
\quad\text{or}\quad
m_j \;\propto\; \bigl|\langle T,[ \gamma_j]\rangle\bigr|
\]
yields a discrete hierarchy indexed by \(w_j\), with no free continuous Yukawa
modulus beyond integers already in the chain.
\end{corollary}

\begin{remark}[Failure Mode]
If \(\mathbf{w}\) is chosen from an open family of windings not constrained by
\eqref{eq:winding-sum} and Section~\ref{sec:closed}, individuality is smuggled
in as a second axiom. Route I is admissible only in the forced form.
\end{remark}

% ---------------------------------------------------------------------
\section{Route II: Leftover \(+1\) on Shells of the 1001-Simplex}
\label{sec:route2}

\begin{definition}[Shells of \(\Delta^{1001}\)]
\label{def:shells}
Let \(\Gamma=\mathbb{Z}_7\times\mathbb{Z}_{11}\times\mathbb{Z}_{13}\) act by the
prime factorization \(1001=7\times 11\times 13\) (equivalently,
\(\Gamma_3=\mathbb{Z}_3\) by the generation qutrit). A \emph{shell} is a
\(\Gamma\)-orbit or a grade of a \(\Gamma\)-equivariant filtration
\[
\Delta^{1001}\;=\;\bigsqcup_{\ell} S_\ell.
\]
Shell indices are integers; there is no continuous radial coordinate.
\end{definition}

\begin{definition}[Generation-Dependent Embedding of \(Q_{\mathrm{res}}\)]
\label{def:embed}
A \emph{residual embedding} assigns the leftover charge \(Q_{\mathrm{res}}=+1\)
to three inequivalent shells:
\begin{equation}
\label{eq:shells}
\iota_{II}\colon g_j \mapsto S_{k_j},\qquad
k_0,k_1,k_2\text{ pairwise distinct}.
\end{equation}
\end{definition}

\begin{lemma}[Single-Shell Obstruction]
\label{lem:single-shell}
If \(Q_{\mathrm{res}}=+1\) is placed in a single shell only, then either
\begin{enumerate}[label=(\alph*)]
  \item the three-generation qutrit collapses to a single effective grade
        (loss of multiplicity), or
  \item flux quantization / entropy conservation across the \(+U/-U\) interface
        fails to match the Global Lemma Core.
\end{enumerate}
Hence any residual embedding compatible with Axiom 0 must occupy (at least)
three inequivalent shells.
\end{lemma}

\begin{theorem}[Route II Individuality]
\label{thm:route2}
Let \(\iota_{II}\) be the canonical residual embedding of
Section~\ref{sec:closed} with shell indices \(\{7,11,13\}\). Then \(\iota_{II}\)
is an admissible individuality structure.
\end{theorem}

\begin{proof}[Proof sketch]
Injectivity is pairwise distinct primes. Discreteness is by
Definition~\ref{def:shells}. The total residual remains \(+1\), so \(G_4\) and
\(|P|\) are unchanged. The identity \(7\times 11\times 13=1001\) and the factor
\(61/1001\) already appear in lifetime formulae.
\end{proof}

\begin{corollary}[Lifetimeage Coupling]
The model lifetime schematic
\[
\tau \;\propto\; \frac{61}{1001}\cdot\frac{\phi^7}{539.9}
\]
admits the generation refinement
\begin{equation}
\label{eq:tau-g}
\tau_j \;\propto\; \frac{61}{1001}\cdot f(k_j)\cdot\frac{\phi^7}{539.9},
\end{equation}
with the fixed rational weight of Section~\ref{sec:closed}
(no free continuous coupling).
\end{corollary}

\begin{remark}[Failure Mode]
Continuous ``shell radius'' reintroduces parameters and is forbidden.
\end{remark}

% ---------------------------------------------------------------------
\section{Route III: Graded Balanced-Ternary Map}
\label{sec:route3}

\begin{definition}[Graded Ternary Data]
\label{def:graded}
A \emph{grading} of the balanced-ternary map is a family of cocycles
\[
c_j\colon\mathbb{N}\to\{-1,0,+1\},\qquad j\in\{0,1,2\},
\]
together with the common underlying map \(T_3\), such that
\begin{equation}
\label{eq:graded-step}
T_3^{(j)}(n)\;:=\;T_3(n)
\quad\text{(set-map unchanged)},
\end{equation}
while observables use weights \(c_j\). The cocycles satisfy
\begin{equation}
\label{eq:cocycle-balance}
\sum_{j=0}^{2} c_j(n)\;=\;Q_{\mathrm{res}}\;=\;+1
\quad\text{for all \(n\) in the master tower}.
\end{equation}
\end{definition}

\begin{definition}[Contraction Rates and Resonant Periods]
\label{def:rates}
Write \(\lambda=\ln 3/539\) for the Banach contraction rate of the ungraded HQCC
map. Graded contractions \(\lambda_j\) must be rational multiples of \(\lambda\).
Graded resonant sub-periods \(\tau_j^{\mathrm{sub}}\) must be drawn from the
sub-harmonic set \(\{5,10,15,30,45\}\,\mathrm{s}\) (or integer commensurators of
\(G_4\))---never three independent fundamental periods.
\end{definition}

\begin{lemma}[Common Orbit Constraint]
\label{lem:common-orbit}
If the three generations are assigned three independent fundamental periods,
fixed-point rigidity (Lemma~\ref{lem:rigidity}) fails. Therefore any admissible
grading keeps a common termination depth \(\sigma=539\pm 1\) and a common
fundamental mode \(G_4\), with individuality in cocycle weights, residue classes
modulo 3, or sub-period structure.
\end{lemma}

\begin{theorem}[Route III Individuality]
\label{thm:route3}
Let \((c_0,c_1,c_2)\) be the canonical graded cocycles of
Section~\ref{sec:closed}. Then
\[
\iota_{III}\colon g_j\mapsto c_j
\]
is an admissible individuality structure: one orbit, one \(G_4\),
generation-dependent dynamical observables.
\end{theorem}

\begin{proof}[Proof sketch]
The set-map \(T_3\) and stopping time \(\sigma\) are generation-independent.
Inequivalence of the \(c_j\) (distinct mod-3 supports) breaks \(\mathrm{Sym}_3\).
Observables (mean sojourn near punctures, weighted transfer spectrum) depend on
\(j\) only through \(c_j\). All numerical inputs are chain integers.
\end{proof}

\begin{remark}[Failure Mode]
Independent fundamental periods not commensurate with \(G_4\) are inadmissible.
\end{remark}

% ---------------------------------------------------------------------
\section{Closed Integer Formulae}
\label{sec:closed}

All ``forced by the chain'' steps are fixed here by explicit integers from
\(\{3,7,11,13,61,539,1001,4880\}\). No continuous choice remains.

\begin{definition}[Canonical Step and Base Winding]
\label{def:Delta}
\[
\Delta \;:=\; |P| \;=\; 61,
\qquad
w_0 \;:=\; \sigma_{\mathrm{HQCC}} \;=\; 539.
\]
\end{definition}

\begin{lemma}[Canonical Windings]
\label{lem:canonical-w}
The unique arithmetic-progression winding assignment with base \(w_0=539\) and
step \(\Delta=61\) is
\begin{equation}
\label{eq:w-explicit}
w_j \;=\; 539 + 61\,j,\qquad j\in\{0,1,2\},
\end{equation}
i.e.
\[
(w_0,w_1,w_2) \;=\; (539,\,600,\,661).
\]
The total class is
\begin{equation}
\label{eq:Wtot}
W_{\mathrm{tot}} \;=\; 539+600+661 \;=\; 1800 \;=\; 3\times 600 \;=\; 3(w_0+\Delta).
\end{equation}
\textbf{Status (No-Go chapter).}
The pair \((w_0,\Delta)=(539,61)\) is \emph{conditional} on a non-circular supply
of \(\sigma=539\) (and of \(|P|=61\)). Flux democracy alone does not select
these anchors (Theorem on flux democracy no-go). By contrast, cocycles
\(c_j\) and shells \(k_j\in\{7,11,13\}\) remain a priori.
\end{lemma}

\begin{proof}
Substitute \eqref{eq:w-explicit} into \eqref{eq:AP-sum}. Summation is elementary.
Both \(539\) and \(61\) are Global Lemma Core integers; no other pair is used.
\end{proof}

\begin{lemma}[Canonical Shells]
\label{lem:canonical-k}
The unique unordered triple of prime shells forced by \(1001=7\times 11\times 13\)
is \(\{7,11,13\}\). Ordering by increasing prime (aligned with increasing
\(j\) and chirality of \(T\)) gives
\begin{equation}
\label{eq:k-explicit}
(k_0,k_1,k_2) \;=\; (7,\,11,\,13).
\end{equation}
\end{lemma}

\begin{definition}[Canonical Dictionaries \(\Phi,\Psi\)]
\label{def:PhiPsi}
\begin{align}
\Phi(w_j) &\;:=\; k_j \;=\; p_{j+1},
  \label{eq:Phi}\\
\Psi(k_j) &\;:=\; c_j,
  \label{eq:Psi}
\end{align}
where \(p_1=7\), \(p_2=11\), \(p_3=13\) are the prime factors of \(1001\) in
increasing order, and \(c_j\) is Definition~\ref{def:c-explicit}. Equivalently,
\[
j \;=\; \frac{w_j-539}{61} \;=\; \mathrm{index}(k_j\text{ in }(7,11,13)),
\]
so \(\Phi\) and \(\Psi\) are determined by the single generation index \(j\).
\end{definition}

\begin{definition}[Canonical Cocycles]
\label{def:c-explicit}
\begin{equation}
\label{eq:c-explicit}
c_j(n) \;:=\;
\begin{cases}
+1 & \text{if } n\equiv j \pmod{3},\\
0 & \text{otherwise}.
\end{cases}
\end{equation}
\end{definition}

\begin{lemma}[Cocycle Balance]
\label{lem:cocycle-ok}
For every \(n\in\mathbb{N}\),
\[
\sum_{j=0}^{2} c_j(n) \;=\; 1 \;=\; Q_{\mathrm{res}}.
\]
Thus \eqref{eq:cocycle-balance} holds pointwise on the master tower.
\end{lemma}

\begin{proof}
Exactly one residue class modulo 3 equals \(j\) for a unique \(j\in\{0,1,2\}\).
\end{proof}

\begin{definition}[Canonical Sub-Harmonics and Contractions]
\label{def:sub-explicit}
\begin{equation}
\label{eq:sub-explicit}
\bigl(\tau_0^{\mathrm{sub}},\tau_1^{\mathrm{sub}},\tau_2^{\mathrm{sub}}\bigr)
\;=\;
(5,\,15,\,45)\,\mathrm{s}
\;\subset\;
\{5,10,15,30,45\}\,\mathrm{s},
\end{equation}
ordered lightest \(\leftrightarrow\) shortest, and
\begin{equation}
\label{eq:lambda-explicit}
\lambda_j \;=\; \lambda \;=\; \frac{\ln 3}{539}
\quad\text{for all }j
\end{equation}
(common contraction; generation dependence lives in \(c_j\) and
\(\tau_j^{\mathrm{sub}}\) only).
\end{definition}

\begin{definition}[Canonical Lifetime Weight]
\label{def:f}
\begin{equation}
\label{eq:f}
f(k) \;:=\; \frac{1001}{k},
\qquad
\bigl(f(7),f(11),f(13)\bigr) \;=\; (143,\,91,\,77).
\end{equation}
\end{definition}

\begin{theorem}[Closure of All Forced Steps]
\label{thm:closure}
With Definitions~\ref{def:Delta}--\ref{def:f} and
Lemmas~\ref{lem:canonical-w}--\ref{lem:cocycle-ok}, every open ``forced by the
chain'' step in Routes I--III is an explicit integer formula. The only free
discrete choice retained is the global orientation
(increasing vs.\ decreasing \(j\)), fixed by the chirality of \(T\) to the
increasing convention \(j=0,1,2\leftrightarrow e,\mu,\tau\).
\end{theorem}

% ---------------------------------------------------------------------
\section{Composite Package}
\label{sec:composite}

\begin{definition}[Composite Individuality Structure]
\label{def:composite}
The \emph{composite package} is
\[
\iota \;=\; (\iota_I,\iota_{II},\iota_{III}),
\]
with
\begin{enumerate}[label=(\roman*)]
  \item Route I: \(g_j\mapsto w_j=539+61j\);
  \item Route II: \(w_j\mapsto k_j\in\{7,11,13\}\) via \(\Phi\);
  \item Route III: \(k_j\mapsto c_j\) via \(\Psi\), on the common \(T_3\) orbit.
\end{enumerate}
Consistency is the commuting square
\begin{equation}
\label{eq:dictionary}
\begin{array}{ccc}
\{g_j\} & \xrightarrow{\;\iota_I\;} & \{w_j\} \\[4pt]
\big\downarrow{\scriptstyle \iota_{III}} & & \big\downarrow{\scriptstyle \Phi} \\[4pt]
\{c_j\} & \xleftarrow{\;\Psi\;} & \{k_j\}
\end{array}
\end{equation}
with \(\Phi,\Psi\) as in Definition~\ref{def:PhiPsi}.
\end{definition}

\begin{theorem}[Individuality without a Second Axiom]
\label{thm:main}
The composite package \(\iota\) is an admissible individuality structure:
\begin{enumerate}[label=(\roman*)]
  \item Axiom 0 remains the sole primitive (cardinality three);
  \item the three generations are pairwise inequivalent as labeled individuals;
  \item \(G_4=539.9\,\mathrm{s}\), HQCC depth \(539\pm 1\), and \(|P|=61\) are
        unchanged;
  \item no free continuous parameter is introduced.
\end{enumerate}
\end{theorem}

\begin{proof}[Proof sketch]
Theorem~\ref{thm:closure} supplies explicit integer labels. Routes I--III are
admissible by Theorems~\ref{thm:route1}--\ref{thm:route3}. The dictionaries
\(\Phi,\Psi\) are bijections on a three-element set, hence introduce no continuum
of moduli. Conditions (i)--(iv) follow.
\end{proof}

\begin{corollary}[Symmetry Breaking Pattern]
\[
\mathrm{Sym}(\{0,1,2\})\;\longrightarrow\;\{1\},
\]
broken completely by \(\iota\). Partial activation of a single route may leave
a \(\mathbb{Z}_2\) stabilizer.
\end{corollary}

% ---------------------------------------------------------------------
\section{Canonical Generation-Label Table (\(e/\mu/\tau\))}
\label{sec:table}

\begin{definition}[Canonical Charged-Lepton Dictionary]
\label{def:leptons}
Identify
\[
(g_0,g_1,g_2) \;=\; (e,\,\mu,\,\tau)
\]
with the closed labels of Section~\ref{sec:closed}. The same \(j\) labels quark
generations when lifted: \((d,s,b)\) and \((u,c,t)\).
\end{definition}

\begin{table}[ht]
\centering
\textbf{Table: Canonical individuality labels for charged leptons}
\label{tab:labels}\\[0.5em]
\begin{tabular}{@{}llrrrrl@{}}
\toprule
Gen. & \(j\) & \(w_j\) & \(k_j\) & \(f(k_j)\) & \(\tau_j^{\mathrm{sub}}\) (s) & \(c_j\) support \\
\midrule
\(e\)   & 0 & 539 & 7  & 143 & 5  & \(n\equiv 0\pmod{3}\) \\
\(\mu\) & 1 & 600 & 11 & 91  & 15 & \(n\equiv 1\pmod{3}\) \\
\(\tau\)& 2 & 661 & 13 & 77  & 45 & \(n\equiv 2\pmod{3}\) \\
\bottomrule
\end{tabular}
\end{table}

\begin{table}[ht]
\centering
\textbf{Table: Global invariants preserved by the dictionary}
\label{tab:invariants}\\[0.5em]
\begin{tabular}{@{}ll@{}}
\toprule
Invariant & Value \\
\midrule
\(W_{\mathrm{tot}}=\sum_j w_j\) & \(1800=3\times 600\) \\
\(\prod_j k_j\) & \(7\times 11\times 13=1001\) \\
\(\sum_j c_j(n)\) & \(1=Q_{\mathrm{res}}=61\bmod 3\) \\
\(\lambda_j\) & \(\ln 3/539\) (common) \\
HQCC depth & \(539\pm 1\) \\
\(G_4\) & \(539.9\,\mathrm{s}\) \\
\(|P|\) & \(61\) \\
\bottomrule
\end{tabular}
\end{table}

\begin{corollary}[Explicit Lifetime Refinement]
Under \eqref{eq:tau-g} and \eqref{eq:f},
\begin{equation}
\label{eq:tau-explicit}
\tau_j \;\propto\; \frac{61}{1001}\cdot\frac{1001}{k_j}\cdot\frac{\phi^7}{539.9}
\;=\; \frac{61}{k_j}\cdot\frac{\phi^7}{539.9},
\end{equation}
so the shell weight cancels one factor of \(1001\) and leaves a pure prime in the
denominator: \(\tau_e\propto 61/7\), \(\tau_\mu\propto 61/11\),
\(\tau_\tau\propto 61/13\), times the common \(\phi^7/539.9\) factor. The standard
muon formula in the theorem chain corresponds to the \(j=1\) line with an overall
normalization \(\tau^-\) and resonance correction \(\delta_{\mathrm{res}}\).
\end{corollary}

\begin{corollary}[Discrete Mass Ordering]
Along Route I, winding gaps are constant:
\[
w_1-w_0 \;=\; w_2-w_1 \;=\; 61,
\]
so any observable monotone in \(w_j\) automatically orders
\(e\prec\mu\prec\tau\) (or the reverse if orientation is flipped---excluded by
chirality of \(T\)).
\end{corollary}

\begin{remark}[Quark Lift]
The same table applies to quarks by reinterpreting \(g_j\) as the \(j\)-th quark
generation. Flavour mixing is then the overlap of cocycles \(c_i,c_j\) on the
common HQCC orbit (CKM / PMNS as discrete cocycle correlators), with no
continuous Yukawa matrix beyond the finite dictionary.
\end{remark}

% ---------------------------------------------------------------------
\section{Comparison and Selection}
\label{sec:compare}

\begin{center}
\begin{tabular}{@{}llll@{}}
\toprule
Route & Carrier & Principal risk & Cleanest fit \\
\midrule
I & Windings on 11D cycles & Free winding family & Brane / flux language \\
II & Shells of \(\Delta^{1001}\) & Continuous radius & \(61/1001\) numerics \\
III & Graded \(T_3\) cocycles & Multiple fundamentals & HQCC core \\
Composite & \((w_j,k_j,c_j)\) & Over-constraining \(\Phi,\Psi\) & Full tower \\
\bottomrule
\end{tabular}
\end{center}

\begin{lemma}[Selection under Rigidity]
\label{lem:selection}
If a single route is preferred, Route III with common orbit and graded cocycle
is the least likely to reopen the fixed point of \(G_4\). If full geometric
language is required, the composite package is preferred, with Route III as the
dynamical readout and Routes I--II as geometric realization. The Axiomatic Book
takes the composite package with Table~\ref{tab:labels} as primary.
\end{lemma}

% ---------------------------------------------------------------------
\section{Falsifiability and Downstream Couplings}
\label{sec:falsify}

\begin{definition}[Individuality Observables]
Generation-dependent observables induced by \(\iota\) include:
\begin{enumerate}[label=(\roman*)]
  \item mass ratios \(m_i/m_j\) as functions of \((w_i-w_j)=61(i-j)\);
  \item lifetime ratios from \eqref{eq:tau-explicit}:
        \(\tau_i/\tau_j = k_j/k_i\) at fixed overall normalization;
  \item mixing angles as overlap of cocycles \(c_i,c_j\) on the common orbit;
  \item sub-harmonic phase offsets at \(\tau_j^{\mathrm{sub}}\in\{5,15,45\}\,\mathrm{s}\).
\end{enumerate}
\end{definition}

\begin{lemma}[Falsifiability]
\label{lem:falsify}
The canonical dictionary is falsifiable if measured mass or lifetime ratios
cannot be matched by Table~\ref{tab:labels} (or its single orientation reverse,
already excluded by chirality) while preserving \(G_4\) and \(|P|=61\). There is
no continuous retuning: the label set is finite.
\end{lemma}

\begin{corollary}[Lifetimeage Ratio Test]
At fixed normalization, \eqref{eq:tau-explicit} predicts
\[
\frac{\tau_\mu}{\tau_e} \;=\; \frac{7}{11},\qquad
\frac{\tau_\tau}{\tau_\mu} \;=\; \frac{11}{13},\qquad
\frac{\tau_\tau}{\tau_e} \;=\; \frac{7}{13},
\]
up to the known overall factors already present in the muon theorem line and
channel-dependent phase space. Gross violation of these discrete ratios (after
those factors) falsifies the canonical shell assignment.
\end{corollary}

% ---------------------------------------------------------------------
\section{Global Lemma Core Form}
\label{sec:glc}

\begin{lemma}[Individuality without a Second Axiom]
\label{lem:glc-individuality}
There exists a discrete, Axiom-0-forced assignment
\[
j\;\mapsto\; (w_j,k_j,c_j)
\;=\;
\bigl(539+61j,\; p_{j+1},\; \mathbf{1}_{n\equiv j\pmod{3}}\bigr)
\]
with \((p_1,p_2,p_3)=(7,11,13)\) such that the three generations are inequivalent
while \(G_4\), HQCC depth, and \(|P|\) remain unique.
\end{lemma}

% ---------------------------------------------------------------------
\section{Summary}
\label{sec:summary}

Axiom 0 yields three as a cardinality. Individuality is restored by the
composite integer dictionary:
\begin{enumerate}
  \item windings \(w_j=539+61j\) on 11-dimensional cycles;
  \item leftover \(+1\) on shells \(k_j\in\{7,11,13\}\) of \(\Delta^{1001}\);
  \item graded cocycles \(c_j=\mathbf{1}_{n\equiv j\pmod{3}}\) on a common \(T_3\)
        orbit of length \(539\pm 1\), with sub-harmonics \((5,15,45)\,\mathrm{s}\).
\end{enumerate}
For charged leptons, \(j=0,1,2\leftrightarrow e,\mu,\tau\)
(Table~\ref{tab:labels}). No second continuous axiom is required; fixed-point
rigidity is preserved; falsifiability is finite and discrete.

% end of chapter
